\documentclass[11pt]{article}

\usepackage[utf8]{inputenc}
\usepackage[T1]{fontenc}
\usepackage{lmodern}
\usepackage{textcomp}
\usepackage[a4paper,margin=1in]{geometry}
\usepackage{xcolor}
\usepackage{graphicx}
\usepackage{booktabs}
\usepackage{array}
\usepackage{longtable}
\usepackage{enumitem}
\usepackage{listings}
\usepackage{fancyhdr}
\usepackage{titlesec}
\usepackage{tcolorbox}
\usepackage{amsmath}
\usepackage{amssymb}
\usepackage[hidelinks]{hyperref}
\usepackage{xurl}

% ---- palette (matches achernar_audit / achernar_commercial) ----
\definecolor{acnavy}{HTML}{1F2A44}
\definecolor{acblue}{HTML}{2D6CDF}
\definecolor{acteal}{HTML}{0E7C7B}
\definecolor{acgray}{HTML}{5B6472}
\definecolor{aclight}{HTML}{EEF2F8}
\definecolor{acred}{HTML}{C0392B}
\definecolor{acamber}{HTML}{B9770E}
\definecolor{acgreen}{HTML}{1E8449}
\definecolor{accode}{HTML}{F4F6FB}

\hypersetup{colorlinks=true,linkcolor=acblue,urlcolor=acblue}

\pagestyle{fancy}
\fancyhf{}
\setlength{\headheight}{14pt}
\addtolength{\topmargin}{-2pt}
\renewcommand{\headrulewidth}{0.4pt}
\lhead{\textcolor{acgray}{\small Achernar — Data Collection \& Ingestion Plan}}
\rhead{\textcolor{acgray}{\small\thepage}}
\lfoot{\textcolor{acgray}{\small Internal — working document}}

\titleformat{\part}{\huge\bfseries\color{acnavy}}{Part \thepart}{0.6em}{}
\titleformat{\section}{\Large\bfseries\color{acnavy}}{\thesection}{0.6em}{}
\titleformat{\subsection}{\large\bfseries\color{acblue}}{\thesubsection}{0.6em}{}
\titleformat{\subsubsection}{\normalsize\bfseries\color{acteal}}{\thesubsubsection}{0.6em}{}

\newtcolorbox{keybox}[1][]{colback=aclight,colframe=acblue,boxrule=0.8pt,
  left=8pt,right=8pt,top=6pt,bottom=6pt,#1}
\newtcolorbox{warnbox}[1][]{colback=white,colframe=acamber,boxrule=0.8pt,
  left=8pt,right=8pt,top=6pt,bottom=6pt,#1}
\newtcolorbox{redbox}[1][]{colback=white,colframe=acred,boxrule=0.8pt,
  left=8pt,right=8pt,top=6pt,bottom=6pt,#1}
\newtcolorbox{okbox}[1][]{colback=white,colframe=acgreen,boxrule=0.8pt,
  left=8pt,right=8pt,top=6pt,bottom=6pt,#1}

\lstset{
  basicstyle=\ttfamily\scriptsize,
  backgroundcolor=\color{accode},
  breaklines=true, frame=single, rulecolor=\color{aclight},
  framesep=6pt, xleftmargin=4pt, xrightmargin=4pt,
  columns=fullflexible, keepspaces=true, showstringspaces=false,
  literate={-}{-}1,
}

% shorthand for code identifiers in prose
\newcommand{\co}[1]{\texttt{\small #1}}

\setlist[itemize]{topsep=2pt,itemsep=1pt,parsep=0pt}
\setlist[enumerate]{topsep=2pt,itemsep=2pt,parsep=0pt}

% long monospaced paths should be allowed to break rather than overflow
\sloppy
\setlength{\emergencystretch}{3em}
\hbadness=10000
\hfuzz=6pt

\begin{document}

\begin{titlepage}
\centering
\vspace*{2.2cm}
{\Huge\bfseries\color{acnavy} Achernar\par}
\vspace{0.5cm}
{\LARGE\color{acblue} Data Collection \& Ingestion Plan\par}
\vspace{0.4cm}
{\large\color{acgray} What has to happen outside, what has to happen in the code,\\
and which of the nine models each night of data actually feeds\par}
\vspace{1.6cm}

\begin{minipage}{0.82\textwidth}
\begin{keybox}
\small
This document is the bridge between \emph{achernar\_v4.pdf} \S7.6--\S7.7 (which
says what the models are) and the repository as it exists on
\today{} (which says what can actually be captured). It is organised so that
Part~I tells you what is broken, Part~II tells you where bytes go, Part~III
tells you what each model needs, and Part~IV tells you exactly what to do at
the eyepiece.

Nothing here assumes a model can be trained on data the current code cannot
record. Where the code cannot record something, that is stated as a blocking
defect with a named file, not as a wish.
\end{keybox}
\end{minipage}

\vfill
{\large\color{acnavy} Ibrahim Shahid\par}
{\small\color{acgray}\today\par}
\end{titlepage}

\tableofcontents
\newpage

% =====================================================================
\part{Instrument and software readiness}
% =====================================================================

\section{Purpose, scope, and how to read this}

\subsection{The question this document answers}

Every model in \S7.6 of the specification is currently a stub, a reference graph
with random weights, or a classical estimator with no fitted constants. The
reason is not that the architectures are unwritten --- most are --- but that
none of them has ever seen a photon collected by this system. The gap between
``the code compiles'' and ``the model ships'' is entirely a data gap.

This document specifies that data: what physically has to be photographed, with
which instrument, under what conditions, how many times, what has to be written
down at the time because it cannot be recovered afterwards, and what the
software has to be able to do before any of it is worth collecting.

\subsection{Reading order}

\begin{itemize}
\item \textbf{Part I} audits the daemon's real capture and sensor-control
      capability, inventories the instruments available, and computes what each
      optical combination can physically resolve. It ends with a list of defects
      that must close before the first useful night.
\item \textbf{Part II} traces the ingestion path that already exists in the
      repository --- \co{capture\_attempts} $\rightarrow$ \co{frames}
      $\rightarrow$ \co{frame\_assets} $\rightarrow$ \co{frame\_measurements}
      --- and identifies where it does not yet reach.
\item \textbf{Part III} takes each of the nine models in turn and states its
      data contract: label source, required covariates, grouping key, sample
      budget, which rig can produce it, and what blocks it today.
\item \textbf{Part IV} is the field material. It is written to the same level of
      detail as the specification chapters: numbered procedures, exact commands,
      frame counts, abort criteria, and the log schema.
\end{itemize}

\subsection{Conventions and standing rules}

\begin{keybox}
\textbf{Rule 0 --- Citation convention.} A reference written ``spec \S7.6.4''
points at \emph{achernar\_v4.pdf}. A bare ``\S7.4'' points at a section of
\emph{this} document. Both documents happen to have a \S7, so the word
\emph{spec} is load-bearing wherever it appears.
\end{keybox}

\begin{keybox}
\textbf{Rule 1 --- Manual first.} Every procedure in Part~IV is executable by
hand. Automation is an overlay, never a prerequisite. If the daemon is down you
should still be able to complete the night and ingest afterwards.

\textbf{Rule 2 --- Unknown is a value.} If a quantity was not measured, it is
recorded as absent, never as a plausible default. The schema already models this
(\co{frames.temperature\_state} $\in$ \{\co{measured}, \co{stale},
\co{missing}\}); the field log must too. A guessed ambient temperature will
silently corrupt the Focus Drift Forecaster and there is no way to detect it
later.

\textbf{Rule 3 --- Raw stays raw.} Nothing that is written to disk during a
session is calibrated, stretched, debayered, or scaled. Calibration constants
travel in the FITS header and the device profile. This is already how
\co{ols\_decode\_raw} behaves and the field procedures must not undo it.

\textbf{Rule 4 --- The group is the run, not the frame.} Splits are by focus run
and by night, never by frame or by star. Two stars from the same frame are
near-duplicates; a random split will produce a model that looks excellent and
generalises to nothing.
\end{keybox}

\newpage
\section{Capability audit: what the daemon can capture today}

This section was written by reading the tree, not the previous handoff notes.
File and line references are to the working tree as of \today.

\subsection{The capture path, end to end}

There are two distinct paths through the camera layer, and they have different
maturity.

\subsubsection*{Path A --- streaming / live view (\co{POST /camera/stream/start})}

\co{OlsCamera::start\_stream} (\co{hardware/camera.rs}) opens a format, pushes
frames into a store, and the vision loop (\co{vision/vision\_loop.rs}) converts
each to a \co{LumaFrame}, estimates background, detects stars, and computes
classical HFR in \co{vision/quality.rs}. This path is the one the autofocus
routine uses. It is the most complete part of the system: 1254 lines of
star-detection and quality code with an explicit saturation gate, an explicit
``gate not performed'' state, and no fabricated scalars.

\subsubsection*{Path B --- still capture (\co{POST /camera/capture})}

\co{OlsCamera::capture\_still} (\co{hardware/camera.rs}, \S~line 671) drives a
DSLR through gphoto2: pause live view, apply and \emph{read back} shutter and
ISO, timestamp, trigger once, drain events to \co{CAPTURE\_COMPLETE}, download
every companion asset, resume. The returned \co{StillCapture} carries every
asset the body produced (native RAW, camera JPEG, and any sidecar), and
\co{storage/capture\_store.rs} commits them with a stage-hash-manifest-fsync-
transaction-rename protocol and a startup \co{reconcile()}.

RAW decoding goes through \co{ols\_decode\_raw} into
\co{hardware/raw\_decode.rs}, which produces \textbf{raw sensor ADU} --- no
black subtraction, no white balance, no rescaling --- and \co{fits.rs} writes
the calibration constants into the header (\co{BUNIT}, \co{BLKLEVEL},
\co{CBLACK0..3}, \co{SATURATE}, \co{XBAYROFF}, \co{YBAYROFF}, \co{BAYERPAT}).

\begin{warnbox}
\textbf{Path B has never been compiled.} The RAW decode work and the
\co{capture\_still} implementation were written in an environment with no Rust
toolchain, no LibRaw and no OpenCV. \co{cargo build} on the Mac is the first
real test of roughly 1,500 lines. Budget a full session for this before any
observing is planned around it. See BP-0 in Part~IV.
\end{warnbox}

\subsection{Driver matrix}

\co{hardware/capabilities.rs} reports, at runtime, which drivers were actually
compiled into the bridge. The three-way distinction it makes --- \emph{compiled
in}, \emph{compiled out}, \emph{exposed in Rust but never wired to the bridge}
--- is exactly the distinction that matters for planning, so it is reproduced
here.

\begin{center}
\scriptsize
\begin{tabular}{@{}llp{0.32\textwidth}p{0.28\textwidth}@{}}
\toprule
\textbf{Subsystem} & \textbf{Driver} & \textbf{Status in code} & \textbf{Relevant to us?} \\
\midrule
Camera & \co{sim}      & always available & yes --- pipeline smoke tests \\
Camera & \co{alpaca}   & always available & only via a bridge host \\
Camera & \co{gphoto2}  & compiled if \co{libgphoto2} found & \textbf{yes --- all 5 DSLR/mirrorless} \\
Camera & \co{uvc}      & compiled if \co{libuvc} found & \textbf{yes --- SV105} \\
Camera & \co{asi}      & compiled if ASI SDK present & only if a loaner has one \\
Camera & \co{svb}      & compiled if SVBony SDK present & \emph{not} the SV105 (see \S\ref{sec:owned}) \\
Camera & \co{toup}     & compiled if ToupTek SDK present & unlikely \\
Camera & \co{indi}     & compiled if \co{libindi} found & fallback for loaner gear \\
Camera & \co{indigo}   & compiled if INDIGO found & unlikely \\
Camera & \co{wdir}     & \textbf{Linux only} & see \S\ref{sec:defects} defect D3 \\
\midrule
Mount  & \co{alpaca}   & always available & loaner GoTo, via bridge \\
Mount  & \co{sim}      & always available & testing \\
Mount  & \co{indi}     & compiled if \co{libindi} found & loaner GoTo, native \\
Mount  & serial/LX200  & \co{hardware/serial\_mount.rs} is a 12-line TODO & \textbf{absent} \\
\midrule
Focuser & \co{alpaca}  & always available & only via a bridge host \\
Focuser & \co{sim}     & always available & testing \\
Focuser & \co{indi}    & compiled if \co{libindi} found & \textbf{the realistic option} \\
Focuser & \co{moonlite}& reported \co{disabled} --- Rust only, no bridge & absent \\
Focuser & USB serial   & reported \co{disabled} --- Rust only, no bridge & absent \\
\midrule
Filter wheel & any     & \co{hardware/filter\_wheel.rs} is \co{todo!()} & \textbf{absent} \\
Env.\ sensors & any    & \co{hardware/sensors.rs} is \co{todo!()} & \textbf{absent} \\
\bottomrule
\end{tabular}
\end{center}

\subsection{What exists above the hardware layer}

It is worth being precise about how much is \emph{not} the problem, because it
changes where effort should go.

\begin{center}
\small
\begin{tabular}{@{}p{0.30\textwidth}p{0.62\textwidth}@{}}
\toprule
\textbf{Area} & \textbf{State} \\
\midrule
Star detection / HFR & Complete and careful. \co{vision/quality.rs}, 1254 lines. Explicit rejection accounting (saturated, edge, blended), an explicit \co{SaturationGate} state distinguishing ``no saturated stars'' from ``never checked''. \\
Background estimation & \co{vision/background.rs} --- implemented, saturation-aware. \\
Trailing (classical) & \co{vision/trailing.rs} --- 223 lines, implemented. \\
FITS writer & \co{fits.rs} --- 507 lines, writes sensor-ADU Bayer frames with full calibration headers. \\
Capture durability & \co{storage/capture\_store.rs} --- staging, hashing, manifest, fsync, transaction, atomic rename, startup reconcile. \\
Schema & 5 migrations, single authority. \co{frame\_assets} models lineage; \co{frame\_measurements} is append-only and versioned by pipeline/model/preprocessing. \\
ONNX serving & \co{vision/fqe\_runner.rs} loads a hash-bound model bundle and publishes shadow-only fields. \co{models/onnx\_runner.rs} generic path is still \co{todo!()}. \\
Debayer / stacking / calibration application & 4-line TODO files. Not needed for training data. \\
Plate solving & \co{processing/plate\_solver.rs} --- 40 lines. \textbf{Needed} (see \S\ref{m:tme}). \\
\bottomrule
\end{tabular}
\end{center}

\subsection{Defects that block data collection}\label{sec:defects}

These are ordered by how much data they spoil, not by how hard they are.

\subsubsection*{D1 --- Sensor saturation point is not declared, and the ML side reads the wrong field}

The daemon models this correctly. \co{vision/quality.rs} distinguishes
\co{max\_adu} (the container ceiling: 255 for an 8-bit stream, 65535 for a
16-bit one) from \co{saturation\_adu} (the sensor's linear saturation point,
\co{Option<f32>}), and \co{resolve\_saturation\_adu} deliberately returns
\co{None} for a 16-bit container rather than 65535, so that the gate reports
``not performed'' instead of silently passing every star.

Two things are still wrong:

\begin{enumerate}
\item \textbf{The value is never set.} \co{[camera] saturation\_adu} is
      commented out in \co{achernar.example.toml}. Until it is measured per body
      and written in, every DSLR frame you take will carry
      \co{saturation\_gate = not\_performed}.
\item \textbf{The dataset builder reads \co{max\_adu}, not \co{saturation\_adu}.}
      \co{ml/data\_pipeline/build\_focus\_quality\_dataset.py} computes its
      saturation level as \co{max\_adu * 0.98} with a default of 64000. On a
      14-bit Canon delivered in a 16-bit container, real saturation is at
      $\sim$16383 ADU; the builder will use 64230 and admit every clipped star
      into the training set. Clipped stars have artificially flat cores and
      therefore artificially \emph{small} measured HFR, so this biases the
      brightest, highest-weight samples toward ``sharper than truth''.
\end{enumerate}

\begin{redbox}
\textbf{Blocking.} Bench procedure BP-2 measures the per-body saturation point
and read noise. The one-line change in the dataset builder --- prefer
\co{saturation\_adu}, and drop frames where it is absent rather than defaulting
--- must land before the first training run. This is the single highest-value
half-hour of code in the document.
\end{redbox}

\subsubsection*{D2 --- \co{saturation\_adu} is plumbed into streaming but not into still capture}

\co{api/camera.rs} passes \co{state.camera\_config.saturation\_adu} into
\co{start\_stream} (line 159). No equivalent call appears on the still-capture
path. That means Path~B frames --- the DSLR RAWs, which are precisely the frames
that need it --- may reach the quality code without it even once D1 is fixed.
Verify during BP-0 and wire it through.

\subsubsection*{D3 --- No focuser with readable absolute position}

The Zhumell Z114 has a manual rack-and-pinion focuser. There is no encoder, no
motor, and no position readout. Consequently:

\begin{itemize}
\item The \textbf{Focus Offset Predictor} (\S7.6.2) cannot be trained --- its
      label \emph{is} a signed step count.
\item The \textbf{Focus Drift Forecaster} (\S7.6.3) cannot be fitted --- $\alpha$
      is $\Delta p / \Delta T$ in steps per $^\circ$C.
\item The \textbf{Focus Quality Estimator} can still be trained (it is
      deliberately pixels-only and takes no position input), but the
      \emph{provenance} of each run is degraded: without positions you cannot
      order samples within a run or verify the sweep was monotonic.
\end{itemize}

Two remedies, in increasing cost, are given in \S\ref{sec:rigs} and RB-2. The cheap one ---
a dial indicator clamped to the drawtube --- turns the manual focuser into a
10\,\textmu m absolute encoder for about \$15 and makes FOP data collectable by
hand.

\subsubsection*{D4 --- No filter wheel, and the hardware layer for one is \co{todo!()}}

\co{hardware/filter\_wheel.rs} is two \co{todo!()} calls. The control layer
above it (\co{filter\_wheel/filter\_planner.rs}, \co{filter\_bandit.rs},
\co{filter\_snr\_scorer.rs}, \co{filter\_safety.rs}) and the schema
(\co{filter\_switch\_events}) both exist. The \textbf{Filter Order Recommender}
(\S7.6.8) is a 7-class model over L/R/G/B/Ha/OIII/SII; you own none of those
filters and no wheel. See \S\ref{m:for} --- the recommendation is to descope, not to
buy.

\subsubsection*{D5 --- No environment sensor}

\co{hardware/sensors.rs} is \co{todo!()}. Temperature currently comes from
\co{read\_temperature()}, i.e.\ whatever the camera reports, with an honest
\co{temperature\_state} when it reports nothing. A DSLR's reported temperature
is the sensor die, which self-heats during a session and is quantised to
$1\,^\circ$C --- useful as a covariate, useless as ambient. FDF needs ambient.
See \S\ref{m:fdf} and RB-4 for the manual-logging workaround, which is adequate for a
first fit.

\subsubsection*{D6 --- Plate solving is a stub}

40 lines in \co{processing/plate\_solver.rs}. This blocks: verified pixel scale
(BP-4), target identity on every frame (needed for the AstroM3 edge retrofit,
\S\ref{m:tme}), airmass and moon-separation covariates computed from a \emph{solved}
pointing rather than a claimed one, and the \co{plate\_solve\_events} table
which already exists. The cheapest path is to shell out to a local
\co{solve-field} (astrometry.net) install rather than implement a solver.

\subsubsection*{D7 --- COC class count disagrees between spec and code}

\S7.6.4 of the specification defines a \emph{binary} classifier:
\co{Linear(1024,1) + Sigmoid}, focal loss with $\alpha=0.25$, $\gamma=2.0$,
temperature scaling on a held-out set. \co{ml/models/cloud\_classifier.py}
declares \co{num\_classes: int = 5} with classes
\{clear, thin, moderate, heavy, overcast\}.

This is not a code detail --- it determines what you write on the log card at
2\,a.m., and relabelling afterwards is impossible because the sky is gone.

\begin{keybox}
\textbf{Resolution adopted in this document:} collect \textbf{5-class} labels in
the field, and derive the binary label by thresholding
(\{clear, thin\} $\rightarrow$ 0, \{moderate, heavy, overcast\} $\rightarrow$ 1).
Five-class labels can always be collapsed to two; two can never be expanded to
five. The model can then be trained either way from the same data. Whichever
head ships, the field protocol is unchanged.
\end{keybox}

\newpage
\section{Instrument inventory and what each rig can actually measure}

\subsection{Owned}\label{sec:owned}

\subsubsection*{SVBony SV105}

\begin{center}
\small
\begin{tabular}{@{}ll@{}}
\toprule
Sensor & Sony IMX307, 1920$\times$1080, 2.9\,\textmu m pixels, 1/2.8" \\
Interface & USB 2.0, \textbf{UVC class} --- appears as a webcam, no vendor SDK \\
Max exposure & \textbf{500\,ms} \\
Output & 8-bit video (UVC YUY2/MJPEG); no 16-bit linear readout \\
Form factor & 1.25" barrel, prime focus in the drawtube \\
\bottomrule
\end{tabular}
\end{center}

Three consequences, all of which shape the plan:

\begin{enumerate}
\item \textbf{It is a UVC device, not an SVBony-SDK device.} The daemon's
      \co{svb} driver (SVBony SDK) will not see it. It must come in through
      \co{uvc} (libuvc). On macOS the OS claims UVC devices for AVFoundation,
      which can prevent libuvc from opening them. \textbf{Verify in BP-1};
      if it fails, the fallbacks are (a) capture with an AVFoundation-based
      tool and ingest via a watch folder --- but \co{wdir} is compiled
      \textbf{Linux-only} --- or (b) run the daemon in a Linux VM/container with
      USB passthrough.
\item \textbf{8-bit output makes the saturation gate load-bearing.} A star peaks
      at 255 and clips. \co{resolve\_saturation\_adu} correctly returns
      \co{Some(max\_adu)} for 8-bit containers, so the gate \emph{will} fire ---
      but only if the frame is tagged as an 8-bit stream. Confirm during BP-1.
\item \textbf{500\,ms ceiling rules out deep sky entirely.} It cannot contribute
      to the exposure bandit, long-exposure FQE, or sky-background-vs-moon
      campaigns. It \emph{is} excellent for high-cadence focus sweeps on bright
      stars, where you want dozens of frames per focuser position in seconds.
\end{enumerate}

\subsubsection*{Zhumell Z114}

114\,mm aperture, 450\,mm focal length, f/3.95, parabolic primary, tabletop
alt-azimuth Dobsonian, 1.25" manual rack-and-pinion focuser, no tracking, no
setting circles.

\begin{keybox}
\textbf{The absence of tracking is an asset, not only a limitation.} An
untracked exposure trails by exactly $15.041\cos\delta$ arcseconds per second.
That is a \emph{physically computed, exactly known} elongation label, available
in unlimited quantity, on any clear night, with no rental. This makes the Z114
the best trailing-detector instrument you have access to --- better than a
tracked loaner rig, which can only produce negatives. See RB-6.
\end{keybox}

\subsection{Loanable camera bodies (UNCG)}

\begin{center}
\small
\begin{tabular}{@{}lllllc@{}}
\toprule
\textbf{Body} & \textbf{Sensor} & \textbf{Pixels} & \textbf{Pitch} & \textbf{RAW} & \textbf{ADC} \\
\midrule
Canon T3i (600D) & APS-C 22.3$\times$14.9 & 5184$\times$3456 & 4.30\,\textmu m & CR2 & 14-bit \\
Canon T6S (760D) & APS-C 22.3$\times$14.9 & 6000$\times$4000 & 3.72\,\textmu m & CR2 & 14-bit \\
Canon M50        & APS-C 22.3$\times$14.9 & 6000$\times$4000 & 3.72\,\textmu m & CR3 & 14-bit \\
Canon R50        & APS-C 22.3$\times$14.8 & 6000$\times$4000 & 3.72\,\textmu m & CR3 & 14-bit \\
Lumix G95        & MFT 17.3$\times$13.0   & 5184$\times$3888 & 3.34\,\textmu m & RW2 & 12-bit \\
\bottomrule
\end{tabular}
\end{center}

\begin{warnbox}
\textbf{The ADC column is a guess until BP-2 measures it, and it must not be
guessed in the config.} The whole point of D1 is that a wrong
\co{saturation\_adu} is worse than none. In particular the G95 is listed as
12-bit here because Panasonic MFT bodies commonly are, but several shoot 12-bit
in electronic shutter and 12-bit in mechanical too; some are 14. Measure it.
\end{warnbox}

\subsubsection*{Tethering risk, by body}

The daemon drives DSLRs through \co{libgphoto2}. Support is per-body and
per-operation, and the operation that matters is \emph{remote shutter release},
not file download.

\begin{center}
\small
\begin{tabular}{@{}lp{0.62\textwidth}@{}}
\toprule
\textbf{Body} & \textbf{Expected gphoto2 status --- to be confirmed in BP-1} \\
\midrule
T3i, T6S & Mature Canon PTP support; remote capture and bulb are the
           best-established case. \textbf{Treat as the primary tethered bodies.} \\
M50      & Listed in the libgphoto2 supported-camera table; users have reported
           intermittent capture failures on the M50 family. Medium risk. \\
R50      & Not on the official supported list. Reports indicate the camera
           connects and live view / file listing work, but
           \co{--capture-image-and-download} may complete without taking a
           frame; behaviour differs between macOS and Linux. \textbf{High risk
           --- do not plan a rental night around it until BP-1 passes.} \\
G95      & Panasonic tether support in gphoto2 is materially weaker than
           Canon's. \textbf{High risk.} Plan for the fallback: shoot with an
           intervalometer, ingest RW2 files afterwards as an offline import. \\
\bottomrule
\end{tabular}
\end{center}

\begin{keybox}
\textbf{Bodies that cannot be tethered are still useful.} A body driven by an
intervalometer produces perfectly good frames for COC (\S\ref{m:coc}), sky background
(\S\ref{m:bkg}) and trailing (\S\ref{m:ted}) --- none of which require the daemon to control the
shutter. What they \emph{cannot} do is participate in a closed loop: no
autofocus sweep, no exposure bandit. Assign untethered bodies to open-loop
campaigns and do not waste rental nights on them.
\end{keybox}

\subsection{The loan-programme telescope}\label{sec:loaner}

The rig arrives shipped, is picked up in Chapel Hill, and is available for a
bounded window. That window is the scarcest resource in this entire plan, so it
should be spent only on the campaigns that \emph{cannot} be run from the
driveway with the Z114. Before booking, get answers to the following --- the
answers change the plan materially, and two of them can make Part~III's Tier~2
disappear entirely.

\begin{enumerate}
\item Aperture, focal length, and focal ratio (fixes pixel scale --- \S\ref{sec:sampling}).
\item Mount model, and \textbf{does it have an ASCOM/Alpaca or INDI driver?}
      The daemon speaks only Alpaca and INDI for mounts. A hand-controller-only
      mount is manual-tracking-only from the daemon's point of view.
\item \textbf{Is there an electronic focuser, and what protocol?} If yes ---
      and it is INDI- or Alpaca-drivable --- then FOP and FDF become
      collectable during the rental and you do not need to buy a focuser. This
      is the single most consequential question on the list.
\item Is there a filter wheel or filter drawer, and which filters?
\item Is a guide camera / guide scope included? (Guiding RMS is a context
      feature for the exposure bandit, \S\ref{m:bandit}.)
\item Loan duration, shipping lead time, and whether the window can be chosen
      to fall on a new moon.
\item Power requirement --- mains, or 12\,V battery? This constrains site
      choice more than people expect.
\end{enumerate}

\subsection{Sampling: what each combination can resolve}\label{sec:sampling}

Plate scale in arcseconds per pixel is $s = 206.265\,p\,/\,f$ with $p$ in
\textmu m and $f$ in mm. Typical Piedmont North Carolina seeing is
$2.5$--$4''$; the table uses $3''$.

\begin{center}
\scriptsize
\begin{longtable}{@{}llrrr@{}}
\toprule
\textbf{Optic ($f$)} & \textbf{Camera} & \textbf{"/px} & \textbf{FOV (arcmin)} & \textbf{FWHM @ 3" (px)} \\
\midrule
\endhead
Z114 native (450\,mm)  & SV105  & 1.33 & 42.5 $\times$ 23.9 & 2.26 \\
                       & T3i    & 1.97 & 170.4 $\times$ 113.6 & 1.52 \\
                       & T6S / M50 / R50 & 1.70 & 170.4 $\times$ 113.6 & 1.76 \\
                       & G95    & 1.53 & 132.2 $\times$ 99.1 & 1.96 \\
\midrule
Z114 + 2$\times$ Barlow (900\,mm) & SV105 & 0.66 & 21.3 $\times$ 12.0 & 4.51 \\
                       & T3i    & 0.99 & 85.2 $\times$ 56.8 & 3.04 \\
                       & T6S / M50 / R50 & 0.85 & 85.2 $\times$ 56.8 & 3.52 \\
                       & G95    & 0.76 & 66.1 $\times$ 49.6 & 3.92 \\
\midrule
80\,mm f/6 refractor (480\,mm) & SV105 & 1.25 & 39.9 $\times$ 22.4 & 2.41 \\
                       & T6S / M50 / R50 & 1.60 & 159.7 $\times$ 106.5 & 1.88 \\
\midrule
8" f/5 Newtonian (1000\,mm) & SV105 & 0.60 & 19.1 $\times$ 10.8 & 5.02 \\
                       & T3i    & 0.89 & 76.7 $\times$ 51.1 & 3.38 \\
                       & T6S / M50 / R50 & 0.77 & 76.7 $\times$ 51.1 & 3.91 \\
                       & G95    & 0.69 & 59.5 $\times$ 44.6 & 4.36 \\
\midrule
8" f/10 SCT (2032\,mm) & SV105  & 0.29 & 9.4 $\times$ 5.3 & 10.19 \\
                       & T6S / M50 / R50 & 0.38 & 37.7 $\times$ 25.2 & 7.95 \\
\midrule
11" f/10 SCT (2800\,mm) & SV105 & 0.21 & 6.8 $\times$ 3.8 & 14.04 \\
                       & T6S / M50 / R50 & 0.27 & 27.4 $\times$ 18.3 & 10.96 \\
\bottomrule
\end{longtable}
\end{center}

\begin{redbox}
\textbf{The Z114 at native focal length is undersampled at best focus for every
camera you have.} Nyquist wants FWHM $\gtrsim 2$ px; the native combinations sit
at 1.5--2.3\,px. The consequence is specific and important: \emph{the labels
near the bottom of the V-curve are quantisation-limited}. The defocused wings
are fine --- they are many pixels wide --- but the region the Focus Quality
Estimator most needs to resolve is exactly the region the optics cannot deliver.

\textbf{Remedy: a 2$\times$ Barlow (roughly \$40--60).} It moves every
combination to 3--4.5\,px FWHM, which is properly sampled, and it doubles the
defocus excursion per millimetre of focuser travel, which makes the dial-
indicator scheme in RB-2 more precise in angular terms. For FQE and TED work on
the Z114 the Barlow should be treated as mandatory equipment, not an option.
\end{redbox}

\begin{warnbox}
\textbf{Unverified: whether a DSLR reaches prime focus on the Z114 at all.}
Many tabletop Newtonians place the focal plane inside the drawtube, so a
camera body with a 44--45\,mm flange distance cannot reach it. The SV105, being
a 1.25" barrel that inserts deep into the tube, almost certainly will. Test this
in daylight before planning anything around it --- BP-3.
\end{warnbox}

\subsection{Rig definitions used throughout this document}\label{sec:rigs}

A ``rig'' is the grouping key that matters for generalisation testing: the
\co{rigs} table exists in migration 002 and \co{focus\_runs.rig\_id} references
it. Hold at least one rig out of training entirely.

\begin{center}
\small
\begin{tabular}{@{}llp{0.42\textwidth}@{}}
\toprule
\textbf{ID} & \textbf{Configuration} & \textbf{What it can produce} \\
\midrule
R1 & Z114 + SV105, native & TED (untracked ladder), fast focus sweeps, planetary/lunar. No deep sky. \\
R2 & Z114 + 2$\times$ Barlow + SV105 & As R1, properly sampled. \textbf{Primary FQE rig without rental.} \\
R3 & Z114 + Barlow + T6S or T3i & Longer exposures, 14-bit, tethered focus sweeps. Untracked, so exposures $\lesssim 1$\,s. \\
R4 & DSLR + kit lens, fixed tripod & COC, sky background, all-sky. No telescope needed. \textbf{Runs in parallel with everything else.} \\
R5 & Loaner scope + loaner mount + DSLR & Tracked long exposure, exposure bandit, FOP/FDF if an electronic focuser is included. \\
\bottomrule
\end{tabular}
\end{center}

\begin{keybox}
R4 is free and independent. It uses a body, a lens, a tripod and an
intervalometer, needs no telescope, no daemon and no rental, and it is the only
rig that collects the Cloud/Obstruction Classifier's data. \textbf{It should be
running on every single night you are outside, including nights whose primary
purpose is something else, and including bad nights.}
\end{keybox}

% =====================================================================
\part{The ingestion pipeline}
% =====================================================================

\section{The path bytes take, as built}

\subsection{Table by table}

The schema is the single authority (\co{achernar-daemon/migrations/}); the Rust
repositories under \co{src/storage/} are its clients. A single exposure touches
four tables.

\begin{center}
\small
\begin{tabular}{@{}lp{0.66\textwidth}@{}}
\toprule
\textbf{Table} & \textbf{Role} \\
\midrule
\co{capture\_attempts} & Written \emph{before} the camera is commanded. Carries
\co{capture\_id}, \co{requested\_json}, and a status of
\co{started|committed|failed}. This is what makes ``the mount was slewing and
the capture silently did nothing'' distinguishable from ``no capture was
requested''. For a data campaign this is the row that lets you count what you
\emph{attempted}, not just what succeeded --- which is the denominator in every
yield calculation in Part~IV. \\
\addlinespace
\co{frames} & One row per exposure. Migration 004 widened it from a thin file
record into a capture record: \co{rig\_id}, \co{camera\_name},
\co{capture\_type} (\co{light|dark|flat|bias}), \co{capture\_status},
\co{exposure\_started\_at} / \co{\_completed\_at}, \co{requested\_exposure\_s},
\co{requested\_gain}, \co{iso}, \co{shutter\_label}, \co{focuser\_position},
\co{temperature\_source}, \co{temperature\_state}, \co{software\_version},
\co{driver\_name}, \co{acquisition\_json}. \\
\addlinespace
\co{frame\_assets} & The bytes. One row per artefact with an explicit
\co{role}: \co{native\_raw}, \co{camera\_jpeg}, \co{sensor\_fits},
\co{analysis\_luma}, \co{thumbnail}, \co{calibrated\_fits}. Carries
\co{checksum\_sha256}, \co{byte\_length}, \co{bit\_depth}, and
\co{derived\_from\_asset\_id} --- so the FITS knows it came from the RAW and the
luma knows it came from the FITS. \co{variant} means a reprocessed asset is a
new row, not an overwrite. Paths are \emph{relative} to the capture root, so the
archive can be moved to an external drive without rewriting the database. \\
\addlinespace
\co{frame\_measurements} & Append-only, one row per (frame, measurement type).
\co{hfr | fwhm | star\_count | fqe | cloud | trailing}, each stamped with
\co{pipeline\_version}, \co{model\_version}, \co{preprocessing\_version}, plus
\co{valid} and \co{rejection\_reason}. Measurements are opinions produced by a
specific version of a specific pipeline; overwriting them inside \co{frames}
would make it impossible to re-evaluate later without destroying what the
previous model said. \\
\bottomrule
\end{tabular}
\end{center}

\subsection{The commit protocol}

\co{storage/capture\_store.rs} does the following, in order, and the order is
the point:

\begin{enumerate}
\item Insert the \co{capture\_attempts} row with status \co{started}.
\item Stage every downloaded asset to \co{<name>.partial}, hashing while
      writing (one pass, not two).
\item Write \co{manifest.v1.json} beside them.
\item \co{fsync} the files and the directory.
\item Open a SQLite transaction; insert \co{frames} + all \co{frame\_assets}.
\item Atomically rename each \co{.partial} to its final name.
\item Commit the transaction; set the attempt to \co{committed}.
\item \emph{Only then} delete the originals from the camera's card.
\end{enumerate}

\co{reconcile()} runs at startup and repairs the two failure modes this leaves:
committed rows whose files are missing, and staged files with no committed row.

\begin{keybox}
\textbf{Field consequence.} If the laptop dies mid-download, nothing is lost and
nothing is half-recorded --- but the frame may still be on the camera card.
\textbf{Do not format cards in the field.} The post-session checklist (Part~IV,
\S D) has a card-retention step for exactly this reason.
\end{keybox}

\subsection{What the ML side reads}

\co{ml/data\_pipeline/build\_focus\_quality\_dataset.py} joins
\co{focus\_samples} to \co{frames} via \co{focus\_samples.frame\_id}, loads the
16-bit PNG at \co{frames.local\_path}, detects stars, and labels each crop with
a SEP/DAOPHOT-style HFR from \co{ml/labels/psf\_measure.py}. It emits an
\co{.npz} keyed by \co{run\_key} (the split/grouping key), plus \co{rig\_id},
\co{camera\_id}, \co{filter\_name}, \co{night\_id} as deployment strata.

\co{ml/data\_pipeline/split\_registry.py} then performs grouped splitting ---
all stars and all frames of one autofocus run stay on the same side --- and
writes a JSON registry so a split is reproducible.

\begin{warnbox}
The builder's own docstring names the wall: ``this needs a populated
\co{achernar.db} and a \co{frames/} directory produced by real autofocus runs on
real rigs. Until that exists the builder has nothing to read.'' Everything in
Part~IV exists to produce exactly that.
\end{warnbox}

\subsection{Gaps between the pipeline and the campaign}\label{sec:gaps}

\begin{center}
\small
\begin{tabular}{@{}lp{0.50\textwidth}l@{}}
\toprule
\textbf{Gap} & \textbf{Description} & \textbf{Blocks} \\
\midrule
G1 & Saturation level read from \co{max\_adu} not \co{saturation\_adu} (D1). & FQE labels \\
G2 & No offline-import path for frames the daemon did not capture. Untethered bodies (G95, possibly R50) and intervalometer runs produce files with no DB row. & R4 entirely \\
G3 & \co{wdir} watch-folder driver is Linux-only. & G2's obvious fix \\
G4 & No ingestion for a manually recorded environment log (ambient temperature, sky condition label, wind). & FDF, COC \\
G5 & \co{capture\_type} exists but nothing sets it to \co{bias|dark|flat}. & calibration library \\
G6 & Plate solve stub, so \co{frames} has no verified pointing, no true pixel scale, no target identity. & TME bookkeeping, airmass covariates \\
\bottomrule
\end{tabular}
\end{center}

\subsubsection*{G2 and G4 are the two that must close first}

They share one solution: a \textbf{sidecar-import tool}. A small Python or Rust
utility that takes (a) a directory of camera files and (b) a session log in the
JSON schema of Part~IV \S C, and writes the corresponding \co{frames} /
\co{frame\_assets} rows using the same hashing and manifest discipline as
\co{capture\_store.rs}.

This is worth building before the first night, because without it every frame
that the daemon did not personally trigger is scientifically orphaned --- and
that is the majority of what the plan collects. It should reuse
\co{capture\_store.rs} rather than reimplement it, and should mark imported rows
with \co{driver\_name = "offline\_import"} and an \co{acquisition\_json}
recording which log file supplied the metadata, so imported and daemon-captured
data can always be separated during analysis.

\section{Grouping, leakage, and the things that cannot be recovered later}

\subsection{The grouping hierarchy}

\begin{center}
\small
\begin{tabular}{@{}llp{0.44\textwidth}@{}}
\toprule
\textbf{Level} & \textbf{Key} & \textbf{Why it is a group} \\
\midrule
Star & --- & Never a split unit. Stars in one frame share seeing, focus, tracking and sky. \\
Frame & \co{frames.id} & Never a split unit. Frames in one run share the optical train and the night. \\
Run & \co{focus\_runs.id} & \textbf{The FQE/FOP split unit.} \\
Night & \co{night\_id} & \textbf{The COC/background/bandit split unit.} Seeing, transparency and sky brightness are night-level. \\
Rig & \co{rigs.id} & Hold one out entirely --- this is the only test of whether the model learned optics or learned physics. \\
\bottomrule
\end{tabular}
\end{center}

\S7.7.2 of the specification additionally requires a \emph{temporal} split
(oldest 70\,\% train / next 15\,\% val / newest 15\,\% test), which the current
\co{split\_registry.create\_group\_splits} does not implement --- it shuffles
groups randomly with a seed. Both are defensible; random-group is a fair test of
\emph{rig} generalisation, temporal is a fair test of \emph{drift}. Report both.

\begin{redbox}
\textbf{Practical consequence for the campaign:} a grouped split with a
temporal component needs \emph{groups}, not \emph{frames}. Twenty thousand star
crops from three focus runs is three data points. This is the single most common
way an astro-ML dataset that looks large turns out to be worthless, and it
drives every sample budget in Part~III.
\end{redbox}

\subsection{Irrecoverable metadata}

The following cannot be reconstructed from the pixels afterwards. If it is not
written down at the time, the frames that depend on it are unusable for the
model that needs it.

\begin{center}
\small
\begin{tabular}{@{}p{0.34\textwidth}p{0.28\textwidth}p{0.30\textwidth}@{}}
\toprule
\textbf{Quantity} & \textbf{Needed by} & \textbf{Source} \\
\midrule
Sky-condition class (5-level) & COC & \textbf{You, at the time} \\
Ambient temperature & FDF, FQE covariate & Probe or hand-logged thermometer \\
Relative humidity, dew point & FOR, bandit, dew diagnosis & Probe or phone weather \\
Focuser absolute position & FOP, FDF & Encoder / dial indicator \\
Which direction the focuser was last moved & FOP (backlash) & Log, or daemon \\
Wind (Beaufort estimate) & TED (distinguishes wind from tracking error) & \textbf{You, at the time} \\
Whether tracking was on & TED & Log \\
Moon altitude and illumination & bandit, FOR, background & Computable from time + site \\
Transparency / naked-eye limiting magnitude & bandit, FOR & \textbf{You, at the time} \\
Dew on the corrector or mirror & everything --- it invalidates the night & \textbf{You, at the time} \\
Anything you touched mid-run & everything & \textbf{You, at the time} \\
\bottomrule
\end{tabular}
\end{center}

\begin{keybox}
The last row is the one that gets forgotten. If you bump the tripod, nudge the
focuser to check something, swap an eyepiece, or a car's headlights sweep the
field --- \textbf{write the timestamp}. One un-annotated disturbance in the
middle of a focus run turns a clean V-curve into a bimodal one, and the fitter
will silently converge on the wrong minimum. There is a dedicated
\emph{Disturbances} block on the field card for this.
\end{keybox}

% =====================================================================
\part{Per-model data requirements}
% =====================================================================

\section{How to read the model sections}

Each subsection below follows the same eight-field template.

\begin{center}
\small
\begin{tabular}{@{}ll@{}}
\toprule
\textbf{Field} & \textbf{Meaning} \\
\midrule
Predicts & The target, in its natural units. \\
Label source & Where ground truth comes from --- and whether a human must supply it. \\
Unit / group & The row of the dataset, and the key splits are made on. \\
Covariates & What must accompany every row, and where it comes from. \\
Rigs & Which of R1--R5 can produce it. \\
Budget & v0 / v1 / v2 sample targets, in \emph{groups} first and rows second. \\
Blocked by & Defects from \S\ref{sec:defects} that must close first. \\
Runbook & The Part~IV procedure that collects it. \\
\bottomrule
\end{tabular}
\end{center}

The three budget tiers mean:

\begin{itemize}
\item \textbf{v0 --- plumbing.} Enough to prove the dataset builder, the split
      registry, the training loop and the ONNX export run end to end on real
      bytes. Not enough to evaluate anything. Expect the model to be bad;
      the deliverable is a green pipeline.
\item \textbf{v1 --- shadow-eligible.} Enough to train something that beats the
      classical baseline on a held-out group split and can run in shadow mode
      alongside the classical authority. Not promotable.
\item \textbf{v2 --- promotion candidate.} Enough groups, rigs, nights and
      covariate span that a held-out \emph{rig} evaluation is meaningful.
\end{itemize}

\section{Model-by-model}

\subsection{Model 1 --- Focus Quality Estimator (FQE), spec \S7.6.1}\label{m:fqe}

\begin{center}
\small
\begin{tabular}{@{}p{0.16\textwidth}p{0.78\textwidth}@{}}
\toprule
Predicts & $\mu = \log(\mathrm{HFR}_{\mathrm{px}})$ and $\log\sigma^2$ per star crop, pooled robustly to a frame value. Note the repo deviates from the spec deliberately: per-star not per-frame, log-HFR in \emph{pixels} not arcseconds, and \textbf{pixels only, no focuser position} (position would let the network regress HFR from a rig-specific shortcut instead of measuring PSF width). \\
Label source & Automatic --- SEP/DAOPHOT HFR from \co{ml/labels/psf\_measure.py}. \textbf{No human labelling.} \\
Unit / group & Row = one star crop. Group = \co{focus\_runs.id}. \\
Covariates & \co{rig\_id}, \co{camera\_id}, \co{night\_id}, \co{filter\_name}, \co{position}, \co{star\_snr}. Position is recorded as provenance, not fed to the net. \\
Rigs & R2 (primary, no rental), R3, R5. \\
Budget & \textbf{v0:} 20 runs, 1 rig, 2 nights. \textbf{v1:} 60 runs, $\geq$2 rigs, $\geq$4 nights. \textbf{v2:} 150 runs, $\geq$3 rigs, $\geq$8 nights, $\geq$10\,$^\circ$C ambient span, one rig fully held out. \\
Blocked by & D1 (saturation), D2, and the Barlow (\S\ref{sec:sampling}). \\
Runbook & RB-2 (manual focuser), RB-3 (motorised). \\
\bottomrule
\end{tabular}
\end{center}

\subsubsection*{Why the budget is expressed in runs}

A single 21-position sweep with 150 usable stars per frame yields
$\sim$3,000 crops --- but they are one group. The binding constraint on a
grouped split is the group count, not the row count. Thirty runs is the
practical floor for a train/val/test split with any power at all; 150 runs
across three rigs is what makes a held-out-rig evaluation interpretable.

At 21 positions $\times$ 5 frames $\times$ $\sim$8\,s per position, one sweep is
roughly 15 minutes including settling. Three sweeps per hour, four productive
hours, gives \textbf{$\sim$10 runs per night} --- so v1 is four to six clear
nights and v2 is fifteen to twenty. That is the real schedule.

\subsubsection*{What makes a run usable}

\begin{itemize}
\item Monotonic on each side of the minimum, with the minimum interior to the
      sweep --- not at an endpoint.
\item $\geq$15 detected stars at every position (fewer at the extremes is
      expected and acceptable down to 8).
\item Saturation gate \emph{performed}, and $<5\,\%$ of stars rejected as
      saturated.
\item No disturbance logged during the run.
\item Ambient temperature drift within the run $<0.5\,^\circ$C, otherwise the
      run is a drift measurement rather than a clean V-curve.
\end{itemize}

\subsection{Model 2 --- Focus Offset Predictor (FOP), spec \S7.6.2}\label{m:fop}

\begin{center}
\small
\begin{tabular}{@{}p{0.16\textwidth}p{0.78\textwidth}@{}}
\toprule
Predicts & Signed focuser steps from the current position to best focus, plus direction logits and an uncertainty. \\
Label source & Automatic, but \emph{retrospective}: the label is $p_{\mathrm{best}} - p_i$, and $p_{\mathrm{best}}$ is only known once the run has converged. \textbf{An aborted run produces no labels at all.} \\
Unit / group & Row = one focus \emph{sample} within a run. Group = run. \\
Covariates & The 38-feature vector of \S7.6.2: 3 quality, 4 classical, 5 position, 4 temperature, 7 filter one-hot, 6 history, 3 direction, 3 rig params, 3 sky. \\
Rigs & R5 with an electronic focuser; R3 with the dial-indicator scheme at reduced precision. \\
Budget & \textbf{v0:} 40 runs. \textbf{v1:} 200 runs ($\sim$4,000 rows). \textbf{v2:} 500 runs ($\sim$10,000 rows), $\geq$3 rigs. \\
Blocked by & \textbf{D3} (no position readout) --- hard blocker. D4 degenerates the 7 filter features to a constant. \\
Runbook & RB-3, and passive harvesting (below). \\
\bottomrule
\end{tabular}
\end{center}

\begin{redbox}
\textbf{FOP needs an order of magnitude more runs than FQE, from the same
nights.} Each run gives FQE thousands of crops but gives FOP only $\sim$20 rows,
one per sampled position, and a LightGBM over 38 features wants $\gtrsim$10,000
rows. Dedicated sweep nights will never get there.

\textbf{The only viable strategy is passive harvesting}: instrument the daemon
so that \emph{every} autofocus run --- during any session, for any purpose ---
writes FOP rows. Then FOP data accumulates as a free byproduct of every other
campaign in this document, and of ordinary use afterwards. The schema for this
already exists (\co{focus\_runs}, \co{focus\_samples}, \co{focus\_actions}); the
work is making sure nothing runs autofocus without logging.
\end{redbox}

\subsubsection*{Feature-group coverage, not just row count}

The 38 features are grouped, and three groups will be degenerate in your data
unless deliberately varied:

\begin{itemize}
\item \textbf{Filter (7 one-hot).} Constant-zero without a wheel. Train with the
      block dropped and document it in \co{norm\_params.json}, or the model
      will carry seven dead inputs into an ONNX graph that a future
      filtered rig will silently misuse.
\item \textbf{Temperature (4).} Needs real thermal span. A single 20-minute
      session has none. Deliberately observe across a night's cooling curve ---
      see RB-4.
\item \textbf{History / direction (9).} \co{prev\_delta\_1/2/3},
      \co{prev\_hfr\_delta\_1/2/3}, \co{dir\_changed\_recently},
      \co{steps\_since\_flip}, \co{consec\_improvements}. These only vary if
      runs actually reverse direction. A sweep that marches monotonically from
      one end to the other never flips. \textbf{Include reversal runs}: start
      inside the V, move the wrong way, detect it, reverse. Roughly one third of
      runs should be reversal runs or the direction head has nothing to learn.
\end{itemize}

The specification's balanced sampler (\S7.6.2, Listing~12) buckets
$|\delta|$ at $[0,5,15,40,120,\infty)$ steps. Sweeps that only ever start near
focus leave the two largest buckets empty and resampling cannot invent them.
\textbf{Deliberately start some runs far out} --- at $\pm$120 steps or more.

\subsection{Model 3 --- Focus Drift Forecaster (FDF), spec \S7.6.3}\label{m:fdf}

\begin{center}
\small
\begin{tabular}{@{}p{0.16\textwidth}p{0.78\textwidth}@{}}
\toprule
Predicts & Focus position at a future time, given temperature; and a refocus trigger when $|p_{\mathrm{forecast}} - p_{\mathrm{now}}| > 0.5\,\mathrm{CFZ}$. \\
Label source & Automatic from consecutive converged runs: $\alpha_{\mathrm{new}} = (p_k - p_{k-1})/(T_k - T_{k-1})$, EWMA'd at $0.8/0.2$ into \co{focuser\_profiles.temp\_coefficient\_steps\_per\_c}. \\
Unit / group & Row = consecutive run pair. Group = night $\times$ rig. \\
Covariates & Ambient temperature (\textbf{not} sensor die temperature), timestamp, rig, focuser position. \\
Rigs & R5, or R3 with dial indicator. \\
Budget & \textbf{v0:} one night, 8 runs, $\geq$4\,$^\circ$C span. \textbf{v1:} 3 nights per rig, $\geq$6\,$^\circ$C each. \textbf{v2:} 3 rigs $\times$ 3 nights, spanning both a cooling and a warming night if possible. \\
Blocked by & \textbf{D3} and \textbf{D5}. \\
Runbook & RB-4. \\
\bottomrule
\end{tabular}
\end{center}

\subsubsection*{This is not really a learned model, and that changes the plan}

FDF v1 is a two-state Kalman filter with hand-set noise parameters
($\sigma_p = 2$ steps, $\sigma_{\dot p} = 0.5$ steps/min, $\sigma_z = 3$ steps)
plus one fitted scalar per rig, $\alpha$. It needs \emph{few, well-separated,
well-thermometered} measurements rather than many noisy ones. Eight clean
autofocus runs spread over a 6\,$^\circ$C night beat two hundred runs at
constant temperature.

Note also that $\alpha$ is a property of the \emph{optical tube}, not of the
camera --- an aluminium tube and a carbon one have different coefficients. A
loaner scope's $\alpha$ does not transfer to the Z114. Each tube needs its own
night.

\begin{keybox}
\textbf{Manual ambient logging is adequate for v1.} A \$10 outdoor
thermometer read every 15 minutes and written on the card, ingested through the
sidecar importer (\S G4), gives $\alpha$ to well within the precision the
Kalman filter's $\sigma_z = 3$ steps assumes. Do not let D5 block this ---
implement \co{hardware/sensors.rs} later, log by hand now.
\end{keybox}

\subsection{Model 4 --- Cloud / Obstruction Classifier (COC), spec \S7.6.4}\label{m:coc}

\begin{center}
\small
\begin{tabular}{@{}p{0.16\textwidth}p{0.78\textwidth}@{}}
\toprule
Predicts & Sky condition. Spec says binary + sigmoid + focal loss; repo stub says 5 classes. \textbf{Collect 5, derive 2} (see D7). \\
Label source & \textbf{Human, at the time of capture.} This is the only model in the system whose labels a person must supply, and they are irrecoverable. \\
Unit / group & Row = one sky frame. Group = \textbf{night}, and within a night, contiguous 10-minute label blocks. \\
Covariates & Timestamp, altitude/azimuth of the field, moon altitude and illumination, whether a light dome is in frame, humidity. \\
Rigs & \textbf{R4} --- body + wide lens + tripod + intervalometer. No telescope, no daemon, no rental. \\
Budget & \textbf{v0:} 2 nights, $\geq$500 labelled frames. \textbf{v1:} 8 nights, $\geq$2,000 frames per class for \{clear, thin\} and $\geq$800 for each of \{moderate, heavy, overcast\}. \textbf{v2:} 20 nights across $\geq$3 moon phases and $\geq$2 sites, with a full held-out night per class. \\
Blocked by & G2 (offline import), G4 (label ingestion). \\
Runbook & RB-5. \\
\bottomrule
\end{tabular}
\end{center}

\begin{keybox}
\textbf{The counterintuitive instruction: go outside when it is cloudy.} Clear
nights are over-represented in any naturally collected set, and the spec already
anticipates this with focal loss ($\alpha = 0.25$, $\gamma = 2.0$). But focal
loss reweights a class that exists; it cannot create one. Broken, marginal,
scudding-cumulus nights are the \emph{most} valuable nights for COC and the
least valuable for everything else --- so they cost nothing to spend this way.

Put R4 on the tripod on nights you would otherwise write off.
\end{keybox}

\subsubsection*{Hard negatives to collect deliberately}

The spec calls for hard-negative mining. These are the ones this site will
actually produce:

\begin{itemize}
\item Thin high cirrus that a photograph barely shows but which kills
      transparency.
\item Moon glow near a bright moon --- \emph{not} cloud, and easily confused
      with it.
\item The Durham / Chapel Hill light dome low in the sky --- a bright gradient
      that is not cloud.
\item Aircraft contrails, both fresh and spread.
\item A tree branch, roof line, or power line intruding into frame --- this is
      \emph{obstruction}, which the model is also named for and which is
      trivially easy to collect and easy to forget.
\item Dew or frost on the lens --- a whole-frame softening that resembles
      overcast. Photograph it on purpose once before wiping it off.
\item Passing car headlights and a neighbour's porch light.
\end{itemize}

\subsubsection*{Labelling discipline}

Label on a fixed cadence (every 10 minutes) rather than per frame, and record
transitions to the minute. Take a phone photo of the sky at each label as an
independent cross-check --- when a label looks wrong during cleaning, the phone
photo is what adjudicates. Photographs of the sky taken with the same body at a
fixed exposure are also directly usable as a second view of the same label.

\subsection{Model 5 --- Trailing / Elongation Detector (TED), spec \S7.6.5}\label{m:ted}

\begin{center}
\small
\begin{tabular}{@{}p{0.16\textwidth}p{0.78\textwidth}@{}}
\toprule
Predicts & Classical: frame ellipticity $e = \tilde\mu_{0.2}\{1 - b_i/a_i\}$, trigger at $e > 0.15$ with $N \geq 10$, plus position angle. CNN: only for the ambiguous band $0.10 < e < 0.20$ or $N < 10$. \\
Label source & \textbf{Physically computed.} An untracked exposure trails by $15.041\cos\delta$ arcsec per second, exactly. \\
Unit / group & Row = one $32\times32$ star cutout. Group = (field, exposure step, night). \\
Covariates & Exposure time, declination, plate scale (from BP-4), tracking on/off, wind estimate, position angle of the trail. \\
Rigs & \textbf{R1/R2 for positives} (Z114 is untracked --- every frame is a labelled trail), R5 for negatives. \\
Budget & \textbf{v0:} 1 night, 3 fields, 12 exposure steps. \textbf{v1:} 3 nights, 5 fields, 12 steps, 20 repeats --- $\sim$3,600 frames, $\gtrsim$100k cutouts. \textbf{v2:} add $\geq$2 tracked nights of negatives and $\geq$1 windy night. \\
Blocked by & BP-4 (need a solved plate scale to convert arcsec to px). \\
Runbook & RB-6. \\
\bottomrule
\end{tabular}
\end{center}

\subsubsection*{The exposure ladder gives a continuous label axis}

At R2 (Z114 + Barlow + SV105, 0.66\,"/px), on the celestial equator:

\begin{center}
\small
\begin{tabular}{@{}rrrl@{}}
\toprule
$t_{\mathrm{exp}}$ (s) & trail (arcsec) & trail (px) & regime \\
\midrule
0.017 (1/60) & 0.26 & 0.4 & indistinguishable from round \\
0.05  & 0.75 & 1.1 & sub-seeing \\
0.10  & 1.50 & 2.3 & marginal --- \textbf{the ambiguous band} \\
0.20  & 3.01 & 4.6 & clearly elongated \\
0.33  & 4.96 & 7.5 & \\
0.50  & 7.52 & 11.4 & \\
\bottomrule
\end{tabular}
\end{center}

\begin{keybox}
The rows in bold are the ones the CNN exists for. The classical estimator
already handles obvious trails; the model is only invoked at
$0.10 < e < 0.20$. \textbf{Sample the ladder densely between 0.05\,s and
0.25\,s} --- that band is where the training signal lives, and a linear sweep of
exposure times will under-sample it.

Declination is the second lever: $\cos\delta$ at $\delta = 70^\circ$ is 0.34, so
a high-declination field at 0.30\,s trails the same as an equatorial field at
0.10\,s. Shooting three declinations $\times$ a dense exposure ladder gives a
well-covered 2-D label grid from a manual telescope in one evening.
\end{keybox}

\subsubsection*{Negatives and confounders}

Perfectly round stars are as necessary as trailed ones, and the Z114 cannot make
them. Sources: (a) the loaner tracked rig, (b) SV105 exposures short enough that
the trail is well below a pixel, at $\leq$0.02\,s. Also collect the confounders
that produce elongation \emph{other} than tracking error --- coma at the field
edge of an f/3.95 Newtonian (which is severe and always present),
astigmatism from a pinched mirror, and wind-shake, which produces a
\emph{random-walk} smear rather than a linear one. The position-angle output
is what distinguishes them: tracking error has a consistent PA across the frame,
coma is radial from the optical axis, wind is incoherent.

\subsection{Model 6 --- Background / SNR Estimator, spec \S7.6.6}\label{m:bkg}

\begin{center}
\small
\begin{tabular}{@{}p{0.16\textwidth}p{0.78\textwidth}@{}}
\toprule
Predicts & \co{background\_adu}, \co{background\_sigma}, \co{snr\_proxy\_median}, \co{saturation\_fraction}. Classical (SEP), always active. Optional DnCNN-lite noise map is v2, GPU-only. \\
Label source & \textbf{No labels --- it needs \emph{constants}}: \co{READNOISE} and \co{SENSOR\_SATURATION} appear literally in the spec's Listing~15. \\
Unit / group & N/A for the classical path. \\
Covariates & Per (body, ISO, temperature) --- gain in $e^-$/ADU, read noise in $e^-$, black level, saturation point. \\
Rigs & \textbf{Bench. Indoors. No sky required.} \\
Budget & 5 bodies $\times$ 4 ISOs $\times$ (20 bias + 12 flat pairs) $\approx$ 1,400 frames. One evening per body. \\
Blocked by & Nothing. \textbf{This can be done tonight, in any weather.} \\
Runbook & BP-2. \\
\bottomrule
\end{tabular}
\end{center}

\begin{okbox}
\textbf{This is the highest-value work available right now.} It requires no
clear sky, no rental, no focuser, and no working daemon. It produces the
constants that D1 needs, which unblocks the FQE labels, which unblocks the
entire focus stack. Every other model in this document consumes its output. If
there is one thing to do before reading further, it is BP-2.
\end{okbox}

The photon-transfer method: for a pair of flats at the same illumination,
\[
g = \frac{(\bar{S}_1 + \bar{S}_2) - (\bar{B}_1 + \bar{B}_2)}
         {\sigma^2_{(S_1 - S_2)} - \sigma^2_{(B_1 - B_2)}}
\quad [e^-/\mathrm{ADU}], \qquad
\sigma_{\mathrm{read}} = g\,\sigma_{\mathrm{bias}} \quad [e^-]
\]
The saturation point is the ADU at which the variance-versus-signal curve turns
over --- \emph{not} the maximum value in the file, which is the container
ceiling. Both go into the device profile; the saturation value additionally goes
into \co{[camera] saturation\_adu}.

\subsection{Model 7 --- Exposure / Gain Recommender (LinUCB), spec \S7.6.7}\label{m:bandit}

\begin{center}
\small
\begin{tabular}{@{}p{0.16\textwidth}p{0.78\textwidth}@{}}
\toprule
Predicts & An arm from $\{10, 30, 60, 90, 120, 180, 300\}$\,s given a 17-dim context. \\
Label source & The reward, computed after the fact from five components: sub accepted ($+1.0$), SNR in range ($+0.5$), sub rejected ($-0.5$), saturation $>5\%$ ($-0.3$), trailing ($-0.2$). \\
Unit / group & Row = one (context, arm, reward) tuple. Group = night. \\
Covariates & Sky background, moon illumination, moon--target separation, filter one-hot, target type one-hot, seeing, airmass, guiding RMS. \\
Rigs & \textbf{R5 only.} 300-second exposures require tracking. The SV105 (500\,ms ceiling) cannot participate at all. \\
Budget & \textbf{v0:} 200 tuples, uniformly randomised. \textbf{v1:} 1,500 tuples, $\geq$6 nights, $\geq$3 moon phases. \textbf{v2:} 5,000 tuples with $\geq$2 target types per session. \\
Blocked by & Rental window; D6 (airmass/moon separation want a solved pointing); no guider means one context feature is permanently absent. \\
Runbook & RB-8. \\
\bottomrule
\end{tabular}
\end{center}

\begin{redbox}
\textbf{Log the propensity or the data is worthless for offline evaluation.}
A contextual bandit is trained online, so it appears to need no dataset. But you
cannot afford to let it explore on live rental nights, which means you need
\emph{offline} policy evaluation --- and every offline estimator (IPS, doubly
robust, self-normalised IPS) requires knowing the probability with which the
logging policy chose each action.

Concretely: run a deliberately randomised logging policy. The first 200 subs
use a uniform policy ($\pi = 1/7$ for every arm); thereafter $\varepsilon$-greedy
with $\varepsilon = 0.2$ recorded per decision. Write the sampled arm
\emph{and} its probability into \co{focus\_predictions}-style provenance.
Without that column, six nights of rental time produce a dataset on which no
unbiased estimate of any alternative policy can be computed.
\end{redbox}

\subsubsection*{Store reward components, not the scalar}

The five reward terms have different measurement provenance --- ``sub accepted''
is a policy decision, ``saturation $>5\%$'' is a measurement, ``trailing'' is
another model's output. Storing only the sum makes it impossible to re-weight the
reward later without re-observing. Store all five in
\co{frame\_measurements.payload\_json} and compute the scalar at training time.

\subsection{Model 8 --- Filter Order Recommender (FOR), spec \S7.6.8}\label{m:for}

\begin{center}
\small
\begin{tabular}{@{}p{0.16\textwidth}p{0.78\textwidth}@{}}
\toprule
Predicts & Next filter, 7-class over L/R/G/B/Ha/OIII/SII, LightGBM multiclass with SHAP explanations. \\
Label source & The filter that was actually chosen, plus its realised SNR. \\
Rigs & \textbf{None available.} \\
Blocked by & \textbf{D4} --- no filter wheel, no filters, and \co{hardware/filter\_wheel.rs} is \co{todo!()}. \\
Runbook & RB-9, documented but descoped. \\
\bottomrule
\end{tabular}
\end{center}

\begin{keybox}
\textbf{Recommendation: descope FOR. Do not buy hardware for it.}

The reasoning is not just cost. FOR sits at the top of the dependency stack ---
its feature set consumes ``historical rejection rate per filter'', ``mean HFR per
filter'', ``sky background per filter'', ``seeing estimate'' and ``transparency
score'', every one of which is an output of a model that does not exist yet.
Building FOR before FQE, background and COC work means training a recommender on
features that are currently constants.

If a loaner rig arrives \emph{with} a wheel, log every filter change through
\co{filter\_switch\_events} as free data and revisit later. If you specifically
want a cheap partial path, a manual 1.25" filter drawer with L + UHC + Ha is
around \$150 and yields a 3-class variant --- but that is a different model from
the one in \S7.6.8 and should be documented as such, not silently substituted.
\end{keybox}

\subsection{Model 9 --- Tri-Modal Embedding (TME / AstroM3), spec \S7.6.9}\label{m:tme}

\begin{center}
\small
\begin{tabular}{@{}p{0.16\textwidth}p{0.78\textwidth}@{}}
\toprule
Predicts & A shared embedding space; used for retrieval, target ranking, anomaly detection. \\
Deployments & \textbf{(a) Cloud/planning:} AstroM3 as published --- light curve + spectrum + metadata. \textbf{Weights already trained and available} (\co{scripts/download\_astrom3.sh} $\rightarrow$ \co{ml/checkpoints/astrom3/}). \textbf{(b) Edge/daemon:} \co{AstroM3Edge} --- a retrofit with a \emph{new image encoder} over image crop + capture metadata + on-mount telemetry, reusing AstroM3's Informer, MetaModel and CLIPLoss. \\
Field data needed & \textbf{Deployment (a): none.} It consumes archival survey data, not your photons. \\
 & \textbf{Deployment (b): the new image encoder is the only untrained part}, and contrastive training from scratch is out of reach at your data scale --- AstroM3 was trained on $\sim$10$^5$ objects. \\
Rigs & Any --- as a byproduct. \\
Blocked by & \textbf{D6} (plate solving). \\
\bottomrule
\end{tabular}
\end{center}

\begin{keybox}
\textbf{TME needs no dedicated observing at all --- it needs bookkeeping.}

Two concrete asks, both cheap:

\begin{enumerate}
\item \textbf{Every frame must carry a resolved target identity.} That means
      plate solving works (D6) and the solved centre is matched to a catalogue
      object. This is free once \co{solve-field} is wired in, and it is
      independently valuable for the bandit's airmass and moon-separation
      features.
\item \textbf{Train the edge image encoder by distillation, not contrastively
      from scratch.} Freeze AstroM3's metadata encoder and its embedding space;
      train only a small image adapter to map your crops into that existing
      space, supervised by the frozen encoders' output for the same resolved
      object. This needs thousands of frames, not $10^5$ objects, and every
      other campaign in this document produces them for free.
\end{enumerate}

Set the threshold at \textbf{$\geq$10,000 plate-solved frames with resolved
target identities} before attempting the adapter. You will cross it during the
FQE and TED campaigns without shooting a single dedicated frame.
\end{keybox}

\section{Priority order and campaign schedule}

\subsection{Dependency tiers}

\begin{center}
\scriptsize
\begin{tabular}{@{}lp{0.30\textwidth}p{0.44\textwidth}@{}}
\toprule
\textbf{Tier} & \textbf{Precondition} & \textbf{Yields} \\
\midrule
\textbf{T0} & None --- indoors, any weather & Background/SNR constants; gphoto2 capability truth; saturation values; focus-reach answer. \textbf{Unblocks D1.} \\
\textbf{T1} & One clear night, own gear + Barlow & TED (RB-6); FQE without positions (RB-2); COC in parallel (RB-5). \\
\textbf{T1b} & Any night, even cloudy & COC alone (RB-5). \textbf{The cheapest data in the plan.} \\
\textbf{T2} & An electronic focuser --- bought, or included with the loaner & FOP, FDF, FQE with full provenance. \\
\textbf{T3} & Rental window & Exposure bandit; tracked negatives for TED; long-exposure FQE; plate-solved frames for TME. \\
\textbf{T4} & Filter wheel & FOR. \textbf{Descoped.} \\
\bottomrule
\end{tabular}
\end{center}

\subsection{The critical path}\label{sec:critical}

\begin{enumerate}
\item \textbf{BP-0} --- compile the daemon on the Mac. Nothing downstream is
      real until this passes. Roughly 1,500 lines of Path~B have never been
      through a compiler.
\item \textbf{BP-2} --- photon transfer curves for all five bodies. Indoors,
      one evening each, weather-independent.
\item \textbf{Fix D1} --- dataset builder reads \co{saturation\_adu}; write the
      measured values into config. Half an hour of code, and it gates every
      focus label.
\item \textbf{BP-1, BP-3} --- gphoto2 truth table and focus-reach test. These
      determine which bodies can be planned around at all.
\item \textbf{Build the sidecar importer} (G2/G4). Without it, R4 --- the
      cheapest and most parallel rig --- produces orphaned files.
\item \textbf{Decide the focuser question} before booking the rental. If the
      loaner includes an INDI/Alpaca focuser, T2 and T3 merge into one window
      and you save both the money and a month. If not, buy one --- FOP and FDF
      are the models the entire \S7 control narrative rests on, and neither is
      collectable without it.
\item \textbf{Start RB-5 immediately} and never stop. It costs nothing, it runs
      on nights that are useless for everything else, and COC needs more
      distinct nights than any other model.
\end{enumerate}

\begin{keybox}
\textbf{Sequencing note on the rental.} The rental window is the scarce
resource, so it should be spent last, not first --- after BP-0/BP-2 have proven
the software path, after the importer exists, and after RB-2/RB-6 have shaken
out the field procedures on gear you can re-shoot for free. Arriving at a
14-day loan with untested capture code is how the window gets spent on
debugging. Aim to book it around a new moon, and pad for shipping.
\end{keybox}

\subsection{Risk register}

\begin{center}
\scriptsize
\begin{tabular}{@{}p{0.30\textwidth}p{0.14\textwidth}p{0.46\textwidth}@{}}
\toprule
\textbf{Risk} & \textbf{Severity} & \textbf{Mitigation} \\
\midrule
Path~B does not compile or crashes on first RAW & High & BP-0 early; keep Path~A (streaming) as the fallback for all FQE work --- it is the mature path. \\
libuvc cannot open the SV105 on macOS & High for R1/R2 & BP-1. Fallback: Linux VM with USB passthrough, or AVFoundation capture + sidecar import. \\
DSLR will not reach prime focus on the Z114 & Medium & BP-3, daytime. Fallback: SV105 only for R1/R2; DSLRs go to R4 and R5. \\
R50 / G95 cannot be tethered & Medium & BP-1. Fallback: assign them to open-loop campaigns (R4) where the shutter is on an intervalometer. \\
Loaner has no electronic focuser & High for FOP/FDF & Ask before booking (\S\ref{sec:loaner}, question 3). Fallback: dial indicator (RB-2) or purchase. \\
Weather --- Piedmont NC summer & High & T1b (COC) converts bad nights into data. T0 is weather-independent. Front-load both. \\
Undersampling at native focal length & Medium & 2$\times$ Barlow, treated as mandatory for R2/R3. \\
Focus runs contaminated by unlogged disturbances & Medium & Disturbance block on the field card; drop the run rather than clean it. \\
\bottomrule
\end{tabular}
\end{center}

% =====================================================================
\part{Where the data comes from}
% =====================================================================

\section{Sourcing routes: the decision frame}

\subsection{The rule that governs this whole part}

\begin{redbox}
\textbf{The own-gear track is the critical path. Everything else is optional
upside, and nothing else may block it.}

Every route in this part except the first has a response latency measured in
weeks and an outcome outside your control --- an email that may not be answered,
a loan programme that may not have a focuser, a rental window that may not fall
on a clear moon. Plan the build and the observing campaign as though none of
them will come through. If one does, it is a bonus track that runs in parallel.

The failure mode this rule exists to prevent is spending September waiting for a
reply instead of printing a bracket.
\end{redbox}

\subsection{What each route can and cannot supply}

Rows are the nine models of Part~III. A tick means the route can produce
training-grade data for that model; a tilde means partial or unverified; a cross
means it structurally cannot.

\begin{center}
\scriptsize
\begin{tabular}{@{}lccccccc@{}}
\toprule
\textbf{Model} & \textbf{Own gear} & \textbf{UNCG/TCO} & \textbf{LCO} & \textbf{Skynet} & \textbf{iTelescope} & \textbf{Seestar} & \textbf{GoTo mount} \\
 & \scriptsize R1--R4 + DIY & \scriptsize free & \scriptsize free & \scriptsize \$65 & \scriptsize \$50/hr & \scriptsize \$349+ & \scriptsize \$200 \\
\midrule
FQE (7.6.1)      & \checkmark\checkmark & $\sim$ & $\times$ & $\times$ & $\times$ & $\sim$ & $\times$ \\
FOP (7.6.2)      & \checkmark\checkmark & \checkmark\checkmark & $\times$ & $\times$ & $\times$ & \checkmark & $\times$ \\
FDF (7.6.3)      & \checkmark\checkmark & \checkmark\checkmark & $\sim$ & $\sim$ & $\sim$ & \checkmark & $\times$ \\
COC (7.6.4)      & \checkmark\checkmark & $\times$ & $\times$ & $\times$ & $\times\times$ & $\times$ & $\times$ \\
TED (7.6.5)      & \checkmark\checkmark & $\sim$ & $\sim$ & $\sim$ & \checkmark\checkmark & \checkmark & $\times$ \\
Background (7.6.6)& \checkmark\checkmark & \checkmark & \checkmark & \checkmark & \checkmark & \checkmark & $\times$ \\
Bandit (7.6.7)   & $\times$ & $\sim$ & $\sim$ & $\sim$ & \checkmark\checkmark & \checkmark & \checkmark \\
FOR (7.6.8)      & $\times$ & $\sim$ & \checkmark & \checkmark & \checkmark\checkmark & $\times$ & $\times$ \\
TME (7.6.9)      & $\sim$ & \checkmark\checkmark & \checkmark & \checkmark & \checkmark & $\sim$ & $\times$ \\
\bottomrule
\end{tabular}
\end{center}

Three things fall out of that table immediately.

\begin{enumerate}
\item \textbf{Own gear is the only route that can produce COC data at all.}
      Every remote and institutional option closes the dome when it is cloudy.
      The Cloud/Obstruction Classifier is structurally uncollectable anywhere
      except a camera you personally point at a bad sky.
\item \textbf{Own gear plus a DIY focuser covers six of nine models.} It is not
      a fallback. It is the main plan.
\item \textbf{The paid routes buy exactly the two models that were already
      descoped} --- the exposure bandit and the filter recommender --- because
      both need hardware you do not have and neither is on the critical path.
\end{enumerate}

\subsection{Switching triggers}

Rather than choosing a route now, define the condition under which each becomes
worth activating.

\begin{center}
\small
\begin{tabular}{@{}p{0.26\textwidth}p{0.66\textwidth}@{}}
\toprule
\textbf{Route} & \textbf{Activate when} \\
\midrule
UNCG / TCO & Immediately --- it costs one email and runs in parallel. \\
LCO archive & Immediately --- public data, no login, no application. \\
LCO Global Sky Clubs & If the Durham club will sponsor an application. \\
Skynet & If \$65 is affordable and you want a scriptable API to write a driver against. \\
Seestar S30/S50 & If the DIY focuser build fails or slips past October, \emph{or} if you specifically want an end-to-end Alpaca integration rig. \\
iTelescope & Only after FQE/FOP/FDF are done and you want FOR and the bandit. Not before. \\
GoTo mount & Never, on this budget. See \S\ref{sec:mounts}. \\
Raleigh loaner & Never. \$25 plus residency proof plus in-person pickup for gear worse than TCO's. \\
\bottomrule
\end{tabular}
\end{center}

\section{The UNCG track}
\label{sec:uncg}

\subsection{Three College Observatory}

TCO is operated jointly by UNCG, NC A\&T, and Guilford College. It holds a
\textbf{0.81\,m (32") Ritchey-Chr\'etien on an equatorial mount} --- the largest
telescope in North Carolina --- at a dark-sky site at 5106 Thompson Mill Road,
Graham, roughly 33 miles and 45 minutes from Greensboro, and about midway
between Greensboro and Chapel Hill.

The instrument that matters is a \textbf{fibre-fed eShel \'echelle
spectrograph}, $R \approx 12{,}000$, covering 3800--7900\,\AA{} continuously.

Public nights run roughly monthly from late September to June; the observatory's
own page notes that summer months are too cloudy to be worth scheduling. That
sets the season boundary for anything involving the site.

\begin{keybox}
\textbf{Prof.\ Anatoly Miroshnichenko is the Director of TCO and of the UNCG
Planetarium}, and has already replied to an initial enquiry saying there is data
to work on. That single fact outranks every paid option in this part.
\end{keybox}

\subsection{What his group actually works on}

Verified from published sources, not inferred:

\begin{itemize}
\item Hot stars with circumstellar material, the B[e] phenomenon, and binary
      systems. Co-author on \emph{Galaxies} \textbf{13}, 67 (2025), affiliated
      to UNCG Physics \& Astronomy.
\item A \textbf{long-term binary spectroscopy programme at the TCO 0.81\,m
      running since 2013}. One recent paper alone uses 71 TCO \'echelle spectra
      spanning 2014--2025; TCO spectroscopy of Polaris\,Aa covers 2015--2024.
\item Photometry used for \textbf{candidate selection}: the 2025 paper builds
      criteria from optical and near-infrared colour indices, finds
      $\sim$40 candidate B[e] objects, then uses spectroscopy to classify them.
\end{itemize}

\subsection{What to ask for, in order}

\subsubsection*{Ask 1 --- the archive project (his benefit)}

The select-then-confirm loop in the 2025 paper is precisely the setting where a
cross-modal embedding helps: hand-built colour cuts are replaced by ranking
against confirmed examples in a learned space over photometry, spectra and
metadata. AstroM3's published weights already exist, so this is an application,
not a training project.

The reason this is a good proposal rather than a speculative one is that it is
\textbf{falsifiable against work already done}: their existing candidate list
and spectroscopic follow-ups form a ready-made validation set. If the method
does not recover them, it does not work.

\subsubsection*{Ask 2 --- the autofocus and engineering logs (your benefit)}

Any observatory running ACP, MaxIm, NINA or TheSkyX writes
(focuser position, HFD/FWHM, temperature) triples every night, plus bias, dark
and flat libraries. These are treated as throwaway operational exhaust. They are
\textbf{exactly the FOP and FDF training set}, and a decade of them may already
exist on disk.

This is the highest-value-per-effort request in the entire document: it costs
the observatory one zip file and could supply, in a single transfer, the
hundreds of focus runs that \S\ref{m:fop} says will otherwise take twenty
nights to accumulate.

Be honest that this serves your project, not theirs. Disguising it as
collaboration is both wrong and transparent.

\subsubsection*{Ask 3 --- engineering time (small, later)}

Not ``may I observe'' but ``may I run a defocus sequence on a bright star during
engineering time.'' Focus sweeps need fifteen minutes and a bright star --- not
dark time, not photometric conditions. Offer explicitly to take the nights
nobody else wants.

\subsection{The second contact}

Dr.\ Alicia Aarnio (UNCG Physics \& Astronomy) works on stellar magnetic fields
and young stars using large datasets and computational methods, and runs an
undergraduate research group. She is the more natural contact for the
software and ML side specifically, and has an existing student pathway.

\subsection{What the UNCG track does and does not unblock}

\begin{center}
\small
\begin{tabular}{@{}p{0.30\textwidth}p{0.62\textwidth}@{}}
\toprule
\textbf{Unblocks} & FOP and FDF, \emph{if} the autofocus logs exist and are shared. TME, via the \'echelle archive. Background/SNR, via their calibration frames. \\
\addlinespace
\textbf{Does not unblock} & COC --- they do not observe through cloud. FQE --- unless engineering time is granted, and even then the runs are theirs, not repeatable on demand. \\
\addlinespace
\textbf{Does not replace} & The DIY focuser. Data you cannot re-collect at will is not a substitute for an instrument you control. \\
\bottomrule
\end{tabular}
\end{center}

\section{Remote networks and purchased hardware}

\subsection{iTelescope}

Commercial remote network. Control is by uploading an ACP plan text file; the
complete directive set is \co{\#FILTER}, \co{\#COUNT}, \co{\#INTERVAL},
\co{\#BINNING}, \co{\#SUBFRAME}, \co{\#DITHER}, \co{\#AUTOFOCUS},
\co{\#AFINTERVAL}, \co{\#FILTEROFFSETS}, \co{\#DEFOCUS}, \co{\#PLATESOLVE},
\co{\#TRACKON}/\co{\#TRACKOFF}, \co{\#DARK}, \co{\#BIAS}, \co{\#EXPRESS}.

\begin{warnbox}
\textbf{There is no directive that commands an absolute focuser position.}
\co{\#DEFOCUS} is documented as ``defocus the image slightly'' with no
parameter --- one binary step. No sweep is possible, so FQE and FOP get nothing.
There is also no ASCOM/INDI/Alpaca surface exposed to the user, so the daemon
cannot drive it at all.
\end{warnbox}

What it \emph{is} good for: \co{\#TRACKON} performs orbital tracking, which
their documentation notes produces ``trailed stars and a stationary minor
planet'' --- deliberate, rate-controlled trailing on a professional rig, which is
excellent TED data. Real L/R/G/B/H$\alpha$/OIII/SII wheels make FOR collectable.
Long tracked exposures with varied moon and airmass make the bandit collectable.
And moon discounts make full-moon time cheap, which is exactly when
sky-brightness data is wanted.

Cost: 1 point $=$ US\$1. Their worked example is 104 points for two hours of
session time on T21 at the Plan-40 tier, i.e.\ \textbf{$\sim$\$52/hour};
rates fall on larger plans. Billing is per-session or per-exposure, with a
typical ``system efficiency'' near 0.7 --- about 42 minutes of exposure per hour
of session.

Education accounts exist, but the application requires an institutional email
address and explicitly rejects public domains such as Gmail. A UNCG address
would be the lever.

\subsection{Skynet (UNC Chapel Hill)}

Roughly twenty optical telescopes across four continents plus the Green Bank
radio telescope, operated from UNC-Chapel Hill. Access is normally restricted to
contributing institutions, but is open to anyone taking the tuition-free,
self-paced \emph{Astronomy with Skynet} courses, for which WebAssign charges
about \$65.

The reason it matters more than iTelescope: \textbf{Skynet has a public,
token-authenticated API} for submitting observations. That is something a daemon
driver could be written against, which makes it Achernar-shaped in a way a web
form is not.

\subsection{Las Cumbres Observatory}

\begin{itemize}
\item \textbf{Global Sky Partners} awards up to 1500 hours per year to
      educational programmes; \textbf{Global Sky Clubs} gives 10 hours on the
      0.4\,m network to school astronomy clubs.
\item The \textbf{science archive is public} --- data without a proprietary
      period is downloadable with no login, and proprietary data becomes public
      after 12 months. Access is by web UI or a JSON API.
\item Pipeline-processed data is available as little as 15 minutes after the
      shutter closes.
\end{itemize}

\begin{okbox}
The LCO public archive is the only route in this part that requires no
application, no payment, no email, and no waiting. It is real FITS with real
headers, available today, and it is enough to exercise the ingestion path, the
background estimator, and the TME bookkeeping while the sky is cloudy.
\end{okbox}

\subsection{Smart telescopes}

The ZWO Seestar S50 (\$499, discontinued early 2026) and S30 (\$349) expose
\textbf{camera, mount and focuser over native Alpaca} in recent firmware; the
open-source \co{seestar\_alp} project provides an alternative controller. An
S50 Pro is expected mid-to-late 2026 at \$599--799.

Because the daemon's \co{alpaca} driver is always compiled and needs no SDK,
a Seestar is a genuine \textbf{drop-in end-to-end integration rig} --- tracking
mount, camera, and absolute-position focuser, all reachable with no new code.

Its limits are optical, not electrical: 50\,mm at f/5 (250\,mm focal length)
gives roughly 2.4\,"/px, badly undersampled, and the optics are fixed. Treat it
as an \emph{integration} instrument and a fallback if the DIY build slips ---
not as an FQE rig.

\subsection{Focusers}

\begin{center}
\small
\begin{tabular}{@{}lrp{0.46\textwidth}@{}}
\toprule
\textbf{Option} & \textbf{Cost} & \textbf{Note} \\
\midrule
Dial indicator & $\sim$\$15 & Manual absolute position to 10\,\textmu m. No automation. \\
\textbf{ESP32 + myFP2ESP32} & $\sim$\$40 & \textbf{Recommended.} Speaks Alpaca, ASCOM, INDI and TCP/IP natively. \\
ZWO EAF & \$179 & \co{indi\_asi\_focuser}. No build risk. \\
\bottomrule
\end{tabular}
\end{center}

\begin{keybox}
\textbf{Why the DIY build must be ESP32 and not Arduino.}
\co{OpenLiveStacker/src/focuser.cpp} is 46 lines and its driver list returns
exactly one entry, \co{\{\{"alpaca:remote","auto"\}\}}. OLS implements no
Moonlite driver, no USB-serial driver and no INDI focuser driver --- despite
\co{hardware/focuser.rs} carrying \co{FocuserDriver} variants that map to
\co{"moonlite"} and \co{"usb"}. Those names point at drivers that do not exist,
which is what \co{capabilities.rs} reports as \co{disabled}.

An Arduino speaking the Moonlite protocol therefore requires writing a new OLS
C++ driver and threading it through the bridge ABI. An ESP32 speaking Alpaca
connects to \co{alpaca\_focuser.cpp}, which is already written and already
compiled, with \textbf{zero new C++ or Rust}.

Secondary benefit: the Alpaca focuser interface carries \co{Temperature} and
\co{TempComp} as standard properties, so a DS18B20 on the same board closes
\textbf{D5} at the same time as \textbf{D3}.
\end{keybox}

\subsection{Mounts --- and why not to buy one}
\label{sec:mounts}

\begin{redbox}
A \$200 GoTo mount solves tracking. \textbf{Tracking is not the blocker.}
Consult the table in \S9.2: a mount unlocks the exposure bandit and nothing
else, and the bandit was already descoped. A \$40 focuser unlocks FOP and FDF,
which are the two models the entire \S7 control narrative rests on.

Sub-\$300 GoTo mounts are additionally alt-azimuth, which introduces field
rotation, and none ship with a drivable focuser. If tracked time is later
needed, TCO has an equatorial 0.81\,m and LCO gives hours away.

\textbf{Spend on the focuser. Do not buy a mount.}
\end{redbox}

\subsection{Standing recommendation}

\begin{enumerate}
\item Build the DIY focuser and observe on own gear. This is the plan.
\item Email UNCG in parallel, expecting nothing on any particular timescale.
\item Pull LCO archive data this week --- it is free and needs no permission.
\item Consider Skynet's \$65 for API access worth writing a driver against.
\item Hold iTelescope in reserve for FOR and the bandit, after the focus stack
      is finished.
\item Buy no mount.
\end{enumerate}

% =====================================================================
\part{Field and bench procedures}
% =====================================================================

\begin{keybox}
Every procedure below states, in order: \textbf{purpose}, \textbf{preconditions},
\textbf{equipment}, \textbf{numbered steps}, \textbf{what is recorded},
\textbf{abort criteria}, and \textbf{acceptance}. A procedure that is aborted
still produces a log entry --- an aborted run is evidence, and the
\co{capture\_attempts} table exists so that failures are counted rather than
vanishing.

\textbf{Nothing here requires the daemon to be working.} Where a step names a
daemon endpoint, a manual equivalent is given.
\end{keybox}

\section*{Appendix A --- Bench procedures (indoors, weather-independent)}
\addcontentsline{toc}{section}{Appendix A --- Bench procedures}

\subsection*{BP-0 --- First compile and ABI verification}
\addcontentsline{toc}{subsection}{BP-0 --- First compile}

\textbf{Purpose.} Determine whether the still-capture and RAW-decode path
(roughly 1,500 lines written but never compiled) actually builds and runs.

\textbf{Preconditions.} Mac with Xcode command line tools, Homebrew, Rust
toolchain.

\begin{enumerate}
\item Install the optional native libraries. Each one flips a capability bit
      that \co{hardware/capabilities.rs} reports at runtime:
\begin{lstlisting}[language=bash]
brew install libraw libgphoto2 libusb opencv libtiff
# optional, enables the INDI camera/mount/focuser drivers:
brew install indi
\end{lstlisting}

\item Build, capturing the log --- the CMake summary lines matter as much as the
      exit code:
\begin{lstlisting}[language=bash]
cd achernar-daemon
cargo build 2>&1 | tee /tmp/build.log
grep -Ei "libraw|gphoto|uvc|indi|Decode:" /tmp/build.log
\end{lstlisting}

\item Expect signature mismatches from the capture work. \textbf{Fix them
      against a clean tree first} and do not tangle them with new features.

\item Run and read the startup banner. It should print the OLS Bridge ABI
      version and whether LibRaw was found.
\begin{lstlisting}[language=bash]
cargo run 2>&1 | head -40
\end{lstlisting}

\item Query the capability report and save it. This file is the ground truth
      for every planning decision in Part~I, and it should be re-saved whenever
      the build environment changes:
\begin{lstlisting}[language=bash]
curl -s localhost:8080/drivers | tee docs/capability-report-$(date +%F).json | python3 -m json.tool
\end{lstlisting}

\item Verify D2: confirm whether \co{camera\_config.saturation\_adu} reaches
      the still-capture path, not only \co{start\_stream}. Search
      \co{src/api/camera.rs} for the capture handler and check what it passes
      into the quality code. If it does not, wire it through now.

\item Run the test suite: \co{cargo test}. Then run the Python side:
      \co{python -m pytest ml/tests/}, which exercises the dataset builder and
      the FQE pipeline against synthetic fixtures.
\end{enumerate}

\textbf{Recorded.} \co{docs/capability-report-YYYY-MM-DD.json}; the build log;
a note of every signature fix made.

\textbf{Abort criteria.} None --- this must be completed.

\textbf{Acceptance.} \co{cargo build} and \co{cargo test} pass; the capability
report shows \co{gphoto2: compiled = true}; the startup banner confirms LibRaw.

\subsection*{BP-1 --- Camera capability probe (per body)}
\addcontentsline{toc}{subsection}{BP-1 --- Camera capability probe}

\textbf{Purpose.} Establish, per body, whether gphoto2 can \emph{trigger} an
exposure --- not merely list files --- and which exposure and ISO values are
actually settable. This determines which bodies can be planned around.

\textbf{Preconditions.} BP-0 passed. \textbf{Turn off the camera's Wi-Fi} and
any companion-app pairing; on Canon bodies set USB connection mode to PC/PTP,
not MTP or ``smartphone''.

\textbf{Time.} 20--30 minutes per body. Repeat for all five.

\begin{enumerate}
\item Connect over USB. On macOS, quit Photos, Image Capture and any
      auto-launching importer first --- they claim the device and gphoto2 will
      report it as busy.

\item Detect and dump abilities:
\begin{lstlisting}[language=bash]
gphoto2 --auto-detect
gphoto2 --abilities
gphoto2 --summary
\end{lstlisting}
      In \co{--abilities}, the decisive line is the capture operations list. If
      it does not include image capture, the body is download-only.

\item Enumerate the settable configuration tree and save it --- this is what
      \co{hardware/dslr\_controls.rs} parses with \co{parse\_shutter\_seconds}
      and \co{parse\_iso}:
\begin{lstlisting}[language=bash]
gphoto2 --list-config > docs/gphoto-config-<body>.txt
gphoto2 --get-config /main/capturesettings/shutterspeed
gphoto2 --get-config /main/imgsettings/iso
gphoto2 --get-config /main/imgsettings/imageformat
\end{lstlisting}
      Record whether \co{bulb} appears among the shutter choices and whether
      \co{/main/actions/eosremoterelease} exists (Canon) --- bulb is required
      for exposures beyond 30\,s.

\item Set RAW-only output, a known shutter and a known ISO, then read them back.
      \emph{Reading back is the point} --- the daemon's capture path applies and
      re-reads for exactly this reason, and a body that silently ignores a set
      is worse than one that refuses it:
\begin{lstlisting}[language=bash]
gphoto2 --set-config imageformat=RAW
gphoto2 --set-config shutterspeed=1/60 --set-config iso=800
gphoto2 --get-config shutterspeed --get-config iso
\end{lstlisting}

\item \textbf{The decisive test.} Attempt three remote captures with download,
      and confirm three \emph{new} files with three \emph{distinct} timestamps:
\begin{lstlisting}[language=bash]
for i in 1 2 3; do
  gphoto2 --capture-image-and-download --filename "probe_${i}.%C"
done
ls -l probe_*
\end{lstlisting}

\item Time the round trip (trigger to file on disk). This sets the realistic
      cadence for every sweep in Appendix~B --- a 6-second round trip turns a
      21-position, 5-frame sweep into 10 minutes of pure transfer.

\item Repeat steps 4--6 with \co{shutterspeed=bulb} if available, using a 5-second
      bulb exposure, and confirm the resulting file's actual exposure in EXIF.

\item Verify LibRaw can decode the result:
\begin{lstlisting}[language=bash]
# any LibRaw sample tool, or:
python3 -c "import rawpy,sys; r=rawpy.imread(sys.argv[1]); print(r.raw_image.shape, r.raw_image.dtype, r.black_level_per_channel, r.white_level, r.raw_pattern)" probe_1.CR3
\end{lstlisting}
      Confirm the reported white level. \textbf{This number is a candidate for
      \co{saturation\_adu}, but do not trust it} --- BP-2 measures it. The
      metadata value is a manufacturer hint and is frequently wrong by a few
      hundred ADU.

\item For the SV105, the equivalent probe is a UVC one:
\begin{lstlisting}[language=bash]
# does the OS expose it, and can libuvc claim it?
system_profiler SPCameraDataType
# then attempt capture through the daemon's uvc driver and check
# the reported stream format is 8-bit
curl -s localhost:8080/camera/formats | python3 -m json.tool
\end{lstlisting}
      Confirm the reported container is 8-bit, so that
      \co{resolve\_saturation\_adu} returns \co{Some(255)} and the saturation
      gate can fire.
\end{enumerate}

\textbf{Recorded.} One row per body in the capability table below, plus the
saved config dumps.

\begin{center}
\small
\begin{tabular}{@{}llllllll@{}}
\toprule
Body & detect & set shutter & readback & capture & bulb & round trip & LibRaw \\
\midrule
T3i & & & & & & & \\
T6S & & & & & & & \\
M50 & & & & & & & \\
R50 & & & & & & & \\
G95 & & & & & & & \\
SV105 & \multicolumn{7}{l}{UVC path: opens / 8-bit confirmed / max exposure confirmed} \\
\bottomrule
\end{tabular}
\end{center}

\textbf{Abort criteria.} If \co{--capture-image-and-download} fails 3 of 3, stop
probing that body and mark it \textbf{open-loop only}. Do not spend more than
30 minutes per body --- the fallback (intervalometer + sidecar import) is
perfectly adequate for COC, background and trailing work.

\textbf{Acceptance.} At least two bodies pass full remote capture. If none do,
the entire closed-loop plan collapses to the SV105 and this must be discovered
now, not on a rental night.

\subsection*{BP-2 --- Photon transfer curve, read noise, and saturation}
\addcontentsline{toc}{subsection}{BP-2 --- Photon transfer curve}

\textbf{Purpose.} Measure, per (body, ISO), the conversion gain $g$ in
$e^-$/ADU, the read noise $\sigma_{\mathrm{read}}$ in $e^-$, the black level,
and the linear saturation point in ADU. These are the constants the
Background/SNR estimator needs (\S\ref{m:bkg}) and the value D1 needs.

\textbf{Preconditions.} BP-1 passed for this body. \textbf{Indoors. Any
weather.}

\textbf{Equipment.} Body; lens or body cap; a uniform light source --- a
white laptop or tablet screen at fixed brightness, viewed slightly defocused,
is more than good enough; a dark room; tripod or stable surface.

\textbf{Time.} 45--60 minutes per body, per ISO. Do at least ISO 100, 400, 1600,
3200.

\begin{enumerate}
\item \textbf{Bias frames.} Body cap on, lens off, in a dark room. Set the
      shortest available shutter (1/4000\,s). Capture \textbf{20 frames}. Do
      not let the sensor heat --- pause 5 seconds between frames.

\item \textbf{Dark frames.} Same, but at each exposure you intend to use in the
      field (1\,s, 10\,s, 30\,s, and bulb 120\,s if applicable), \textbf{10
      frames each}. These also give you the dark-current rate, which is
      strongly temperature dependent --- record the EXIF sensor temperature for
      every frame.

\item \textbf{Flat pairs.} Point at the uniform source. At each of
      \textbf{12 illumination levels} spanning roughly 5\,\% to 105\,\% of full
      scale (vary shutter speed, not screen brightness --- brightness controls
      on screens are often non-linear and PWM-modulated), capture
      \textbf{2 frames}. Twelve levels $\times$ 2 = 24 frames.

      Use exposures $\geq$1/100\,s to avoid screen PWM banding, and check one
      frame for banding before committing to the whole ladder.

\item \textbf{Compute.} For each level, using the central 500$\times$500 region
      of one Bayer channel only (do \emph{not} mix channels --- they have
      different sensitivities):
      \begin{itemize}
      \item mean signal $\bar{S}$ and the variance of the \emph{difference}
            of the pair, $\sigma^2_{(S_1-S_2)}$;
      \item the same for the bias pair;
      \item $g = \big[(\bar{S}_1 + \bar{S}_2) - (\bar{B}_1 + \bar{B}_2)\big] \big/ \big[\sigma^2_{(S_1-S_2)} - \sigma^2_{(B_1-B_2)}\big]$.
      \end{itemize}
      Differencing is what removes fixed-pattern non-uniformity, which would
      otherwise dominate the variance and make $g$ come out far too small.

\item \textbf{Plot} $\sigma^2$ against $\bar{S}$. The linear region has slope
      $1/g$. \textbf{The saturation point is where it turns over}, not the
      maximum value in the file.

\item \textbf{Read noise.} $\sigma_{\mathrm{read}} = g \times \sigma_{\mathrm{bias}}$,
      where $\sigma_{\mathrm{bias}}$ is the standard deviation of the bias
      difference frame divided by $\sqrt{2}$.

\item \textbf{Black level.} The mean of the bias frames. Canon bodies typically
      carry an offset near 2048 ADU for 14-bit output. \textbf{Measure it, do
      not assume it} --- it varies by body and sometimes by ISO.

\item \textbf{Write the results} to \co{docs/device-profiles/<body>.json} and
      set \co{[camera] saturation\_adu} in the daemon config for the body in
      use. Follow the config file's own warning: a wrong value is worse than
      none.
\end{enumerate}

\textbf{Recorded.} Per (body, ISO): $g$, $\sigma_{\mathrm{read}}$, black level,
saturation ADU, full-well in $e^-$, dark current in $e^-$/px/s at the measured
sensor temperature.

\textbf{Abort criteria.} If the PTC is not linear over at least a decade of
signal, something is applying a non-linear transform --- almost always a
highlight-tone curve or the wrong RAW decode path. Stop and fix; the numbers are
meaningless otherwise.

\textbf{Acceptance.} $g$ within a factor of 2 of the manufacturer's published
value where one exists; $\sigma^2$-vs-$\bar{S}$ linear ($R^2 > 0.99$) below
saturation.

\subsection*{BP-3 --- Prime-focus reach and backfocus (daytime)}
\addcontentsline{toc}{subsection}{BP-3 --- Prime-focus reach}

\textbf{Purpose.} Determine whether each camera can physically reach focus on
the Z114, with and without the Barlow. Many tabletop Newtonians place the focal
plane inside the drawtube, which no DSLR can reach.

\textbf{Time.} 15 minutes. \textbf{Daylight. Do not point anywhere near the
sun.}

\begin{enumerate}
\item Set up indoors, aimed through a window at a high-contrast object at least
      1\,km away --- a distant roofline or antenna.
\item Insert the SV105 at prime focus. Rack the focuser fully in, then out
      slowly, watching the live stream. Record whether focus is achieved and,
      if so, the drawtube extension in mm at focus.
\item Repeat with each DSLR on a T-ring + 1.25" nosepiece.
\item Repeat both with the 2$\times$ Barlow inserted --- a Barlow moves the
      focal plane \emph{outward}, which usually rescues a body that could not
      otherwise reach focus. This is a second, independent reason to own one.
\item If a DSLR cannot reach focus even with the Barlow, record the shortfall in
      mm. Remedies, in order of preference: use the Barlow permanently; fit a
      low-profile focuser; move the primary mirror up the tube (last resort ---
      it changes the optical alignment).
\end{enumerate}

\textbf{Recorded.} A 6-row table: camera $\times$ \{native, Barlow\}
$\rightarrow$ \{reaches focus?, drawtube mm at focus, shortfall mm\}.

\textbf{Acceptance.} At least one camera reaches focus in at least one
configuration. If none do, R1--R3 are void and every telescope-based campaign
depends on the rental.

\subsection*{BP-4 --- Plate scale and orientation by solving}
\addcontentsline{toc}{subsection}{BP-4 --- Plate scale by solving}

\textbf{Purpose.} Measure the true arcsec/px and field rotation. \textbf{Do not
trust the nominal focal length} --- a Barlow's true amplification depends on its
spacing from the sensor and is routinely 10--20\,\% off its label, and the trail
lengths in RB-6 are computed in arcseconds but labelled in pixels.

\textbf{Preconditions.} One clear night, 20 minutes. Astrometry.net installed
locally with the appropriate index files (4200-series for wide fields,
4100-series for narrow).

\begin{enumerate}
\item Point at a moderately rich star field well away from the Milky Way plane
      --- too many stars is as bad as too few.
\item Capture 3 frames at the shortest exposure that shows $\geq$30 stars.
\item Solve:
\begin{lstlisting}[language=bash]
solve-field --scale-units arcsecperpix --scale-low 0.2 --scale-high 3.0 \
            --downsample 2 --no-plots --overwrite frame.fits
\end{lstlisting}
\item Record from the WCS: pixel scale, field rotation, and the solved centre.
\item Repeat for every optical configuration you intend to use (native, Barlow,
      each camera). \textbf{Each combination has its own scale.}
\end{enumerate}

\textbf{Recorded.} \co{docs/plate-scales.json}, keyed by rig ID.

\textbf{Acceptance.} Solved scale within 15\,\% of the computed value in
\S\ref{sec:sampling}. A larger discrepancy means the assumed focal length or pixel pitch is
wrong --- find out which before collecting anything that depends on it.

\subsection*{BP-5 --- Focuser step calibration (only if motorised)}
\addcontentsline{toc}{subsection}{BP-5 --- Focuser step calibration}

\textbf{Purpose.} Populate \co{focuser\_profiles.steps\_per\_micron} (migration
005) and the critical focus zone, both of which the FDF trigger
($0.5 \times \mathrm{CFZ}$) and the FOP normalisation depend on.

\begin{enumerate}
\item Command a known large move (e.g.\ 1000 steps) and measure the drawtube
      displacement with a dial indicator. $\mathrm{steps/\mu m} = 1000 / \Delta x_{\mu m}$.
\item Measure backlash: approach a position from each direction and record the
      difference. Write it to \co{focuser\_profiles}; \co{focus/backlash.rs}
      consumes it.
\item Compute the CFZ from the focal ratio: $\mathrm{CFZ} \approx 4.88\,\lambda N^2$
      in \textmu m, with $\lambda = 0.55\,\mathrm{\mu m}$ and $N$ the focal ratio.
      For f/3.95 this is $\approx 42\,\mathrm{\mu m}$; for f/10, $\approx 268\,\mathrm{\mu m}$.
      Convert to steps with the measured scale.
\item \textbf{Sanity check.} The step size should be smaller than CFZ/5, or the
      focuser cannot resolve best focus regardless of how good the model is.
\end{enumerate}

\textbf{Recorded.} steps/\textmu m, backlash in steps, CFZ in \textmu m and in
steps, per rig.

\section*{Appendix B --- Observing runbooks}
\addcontentsline{toc}{section}{Appendix B --- Observing runbooks}

\subsection*{RB-1 --- Calibration frame library}
\addcontentsline{toc}{subsection}{RB-1 --- Calibration library}

\textbf{Purpose.} Bias, dark and flat libraries per (body, ISO, exposure,
temperature bin). Required for background estimation and any later calibration;
also the substrate for BP-2.

\textbf{When.} Refresh whenever the body changes, the ISO set changes, or
ambient temperature moves by more than 5\,$^\circ$C from the library's bin.

\begin{center}
\small
\begin{tabular}{@{}llll@{}}
\toprule
\textbf{Type} & \textbf{Count} & \textbf{Conditions} & \co{capture\_type} \\
\midrule
Bias & 50 & shortest shutter, cap on, dark room & \co{bias} \\
Dark & 20 per exposure & matched exposure, ISO, temperature bin & \co{dark} \\
Flat & 20 per filter/config & even illumination, $\sim$50\,\% full scale & \co{flat} \\
Flat-dark & 20 & matched to the flat exposure & \co{dark} \\
\bottomrule
\end{tabular}
\end{center}

\begin{enumerate}
\item Shoot bias and dark at the end of the session, while the equipment is
      still at ambient. Darks shot warm in the afternoon do not match lights
      shot cold at 2\,a.m.
\item Shoot flats at dusk against a twilight sky, or with a white t-shirt over
      the aperture and an even light source. Target a mean around half of the
      \emph{measured} saturation from BP-2 --- not half the container ceiling.
\item Record the sensor temperature (EXIF) for every calibration frame.
\item Set \co{frames.capture\_type} correctly on ingest (gap G5). If the daemon
      does not yet set it, the sidecar importer must.
\end{enumerate}

\textbf{Abort criteria.} Any light leak visible in a bias frame --- reshoot in a
darker room, and check the viewfinder eyepiece is covered on DSLRs.

\textbf{Acceptance.} Bias frames have a flat histogram centred on the black
level from BP-2 with no gradient; flats vary by $<20\,\%$ corner to centre
beyond expected vignetting.

\subsection*{RB-2 --- Focus sweep with a manual focuser (R2 / R3)}
\addcontentsline{toc}{subsection}{RB-2 --- Manual focus sweep}

\textbf{Purpose.} Produce FQE training runs --- and, with the dial-indicator
modification, low-precision FOP runs --- without a motorised focuser.

\textbf{Preconditions.} BP-2, BP-3, BP-4 complete. Clear sky, seeing not
terrible, wind below about 10\,km/h.

\textbf{Equipment.} Z114 + 2$\times$ Barlow + SV105 (or DSLR). A dial indicator
(0.01\,mm) clamped so its plunger rides the drawtube, or failing that a printed
paper mm scale taped along the drawtube. Red headlamp. Field card.

\begin{keybox}
\textbf{The dial-indicator modification.} A \$15 dial indicator clamped to the
focuser body with its plunger against the drawtube gives absolute position to
10\,\textmu m. At f/3.95 the CFZ is $\approx 42\,\mathrm{\mu m}$, so 10\,\textmu m
resolves the focus zone into four steps --- coarse, but \emph{real absolute
position}, which is the difference between an FQE-only run and an FOP-eligible
one. Record the indicator reading in \co{frames.focuser\_position} in units of
10\,\textmu m and document the unit in \co{acquisition\_json}.
\end{keybox}

\begin{enumerate}
\item \textbf{Setup.} Level the table. Let the telescope reach ambient for
      $\geq$30 minutes --- a warm mirror produces tube currents that ruin every
      HFR measurement. Record the start time and ambient temperature.

\item \textbf{Target.} A field with 20--60 stars of moderate brightness. Avoid
      a single very bright star: it saturates while the field stars are still
      undetected, and the saturation gate will reject exactly the star you were
      aiming at. \textbf{A loose open cluster is ideal.}

\item \textbf{Exposure.} Set so the brightest stars peak at 50--70\,\% of the
      measured saturation \emph{when in focus} --- because they will be far
      dimmer when defocused, and you must not clip at the centre of the sweep.
      For the SV105 this will be in the 50--300\,ms range.

\item \textbf{Find visual best focus} and record the indicator reading as
      $x_0$.

\item \textbf{Define the sweep.} 21 positions, symmetric about $x_0$, step size
      $\approx \mathrm{CFZ}/3 \approx 14\,\mathrm{\mu m}$ --- so
      $x_0 \pm 140\,\mathrm{\mu m}$ end to end. If the V-curve does not clearly
      turn over at both ends, widen the range and redo.

\item \textbf{At each position, in order:}
      \begin{enumerate}[label=\alph*.]
      \item Move the focuser \textbf{always in the same direction} --- start
            fully at one end and march to the other. Never back up mid-sweep;
            backlash will corrupt the position axis.
      \item Wait 3 seconds for vibration to settle.
      \item Record the indicator reading.
      \item Capture \textbf{5 frames}.
      \item Note star count from the live display if available.
      \end{enumerate}

\item \textbf{Log the run} on the field card: run ID, start/end time, $x_0$,
      step size, direction, ambient temperature at start and end, target,
      altitude, and anything you touched.

\item \textbf{Repeat} 3 times per night minimum. Vary the target field between
      runs --- three sweeps on the same field are more correlated than three on
      different fields, and correlated groups do not help a grouped split.

\item \textbf{One run in three should be a reversal run} (see \S\ref{m:fop}): start
      70\,\textmu m inside the V on the wrong side, move away from focus for
      four positions, then reverse and complete the sweep. This is the only
      source of variation for the FOP direction and history features.
\end{enumerate}

\textbf{Recorded.} Per run: 21 positions $\times$ 5 frames $\approx$ 105 frames.
Per position: indicator reading, timestamp, exposure, gain/ISO, star count.
Per run: ambient temperature at start and end, target, altitude, seeing
impression, disturbances.

\textbf{Abort criteria.}
\begin{itemize}
\item Cloud crossing the field --- abort, log the run as \co{aborted}, restart.
\item Star count $<8$ at the centre of the sweep --- the field is too sparse.
\item Ambient temperature drifts $>0.5\,^\circ$C within the run --- keep the
      data but flag it; it belongs to RB-4, not here.
\item Wind gusts visibly shaking the image --- abort.
\end{itemize}

\textbf{Acceptance.} The HFR-vs-position curve is a clean V or U, monotonic on
each side, minimum interior to the sweep, with $\geq$15 stars at the minimum and
$\geq$8 at both ends.

\textbf{Yield.} $\sim$15 min per run, $\sim$10 runs per productive night.
\textbf{v1 (60 runs) is 6 clear nights; v2 (150 runs) is 15.}

\subsection*{RB-3 --- Focus sweep with a motorised focuser (R5)}
\addcontentsline{toc}{subsection}{RB-3 --- Motorised focus sweep}

\textbf{Purpose.} Full-provenance FQE runs plus FOP rows and FDF pairs. This is
the procedure that produces all three focus models at once and is therefore the
best use of a rental night.

\textbf{Preconditions.} BP-5 complete; focuser reachable through INDI or Alpaca;
\co{POST /focus/run} works.

\begin{enumerate}
\item Setup and thermal equalisation as RB-2 steps 1--3.
\item Verify the daemon can read and command position:
\begin{lstlisting}[language=bash]
curl -s localhost:8080/focuser/position
curl -s -X POST localhost:8080/focuser/moveto -d '{"position": 12000}'
\end{lstlisting}
\item \textbf{25 positions}, step size $\mathrm{CFZ}/4$ in steps from BP-5,
      \textbf{3 frames} per position, always approaching from the same
      direction with backlash compensation enabled.
\item After each run, record the fitted best position and the ambient
      temperature --- this pair is one FDF observation.
\item \textbf{Vary deliberately across runs:}
      \begin{itemize}
      \item \textbf{Field position:} centre, and two fields near the edge of the
            usable field. Off-axis aberration changes the PSF shape and the
            model must not learn that the centre crop is the only valid one.
      \item \textbf{Airmass:} one field near zenith, one at $\sim$2 airmasses.
      \item \textbf{Starting offset:} alternate between starting near focus and
            starting at $\pm$150 steps, to populate the balanced-sampler buckets
            of \S7.6.2.
      \item \textbf{Direction:} alternate inward-first and outward-first.
      \end{itemize}
\item Log every run to \co{focus\_runs} with a correct \co{rig\_id}.
\end{enumerate}

\textbf{Yield.} $\sim$8 min per run with a motor. \textbf{$\sim$25 runs per
night}, which is where the FOP budget becomes achievable --- 200 runs is 8
nights rather than 20.

\subsection*{RB-4 --- Thermal drift night (FDF)}
\addcontentsline{toc}{subsection}{RB-4 --- Thermal drift night}

\textbf{Purpose.} Fit the per-tube temperature coefficient
$\alpha$ in steps/$^\circ$C, and the drift rate $\dot p$.

\textbf{Preconditions.} A focuser with position readout. A thermometer.
\textbf{Choose a night with a forecast temperature drop of $\geq$8\,$^\circ$C}
--- clear, dry, post-frontal nights are ideal, which conveniently are also good
seeing nights.

\begin{enumerate}
\item Start \textbf{at dusk}, while the air is still warm, and run until the
      pre-dawn minimum. The span is the whole point; a 3-hour session in the
      middle of the night has almost none.
\item Do \textbf{not} let the telescope reach thermal equilibrium first --- for
      once, you want the cooling transient.
\item Run a full autofocus \textbf{every 30 minutes}, on the same field if
      possible, and \textbf{do not touch the focuser between runs}.
\item Record at every run: ambient temperature, sensor temperature, focuser
      best position, time.
\item Additionally log ambient temperature every 15 minutes regardless, so the
      temperature series is denser than the focus series.
\item Continue for $\geq$6 hours, giving $\geq$12 runs.
\end{enumerate}

\textbf{Recorded.} A table of (time, $T_{\mathrm{ambient}}$,
$T_{\mathrm{sensor}}$, $p_{\mathrm{best}}$, run\_id). This is the entire FDF
dataset for that tube.

\textbf{Abort criteria.} Cloud for more than 30 minutes --- resume and note the
gap; do not interpolate across it. If the focuser is bumped, the night is over
for FDF purposes (though the individual runs remain valid FQE data).

\textbf{Acceptance.} $\geq$10 successful runs spanning $\geq$6\,$^\circ$C,
and $p_{\mathrm{best}}$ versus $T$ approximately linear. A scatter plot with no
visible trend means either the temperature range was too small or the focuser
has more backlash than drift --- check BP-5.

\textbf{Repeat.} Three nights per tube. $\alpha$ is a property of the tube, not
the camera; a loaner's $\alpha$ does not transfer to the Z114.

\subsection*{RB-5 --- Cloud and obstruction campaign (R4)}
\addcontentsline{toc}{subsection}{RB-5 --- Cloud campaign}

\textbf{Purpose.} COC training data. \textbf{The cheapest and most parallelisable
campaign in the plan.}

\textbf{Preconditions.} None. No telescope, no daemon, no rental, no clear sky.

\textbf{Equipment.} Any body + widest lens; tripod; intervalometer or the
body's built-in interval timer; spare battery (this runs all night); a phone for
cross-check photos; the field card.

\begin{enumerate}
\item Set up pointing at a fixed altitude/azimuth --- roughly 45$^\circ$
      altitude, and choose the azimuth deliberately. Run some nights aimed at
      the Durham/Chapel Hill light dome and some aimed away, so the model sees
      both. Record which.
\item Manual mode, fixed exposure, fixed ISO, fixed focus at infinity, fixed
      white balance. \textbf{Nothing may auto-adjust} --- an auto-exposure
      change looks exactly like a transparency change and will teach the model
      the wrong thing. RAW + JPEG.
\item Interval: \textbf{one frame every 30 seconds}, all night. At 8 hours
      that is 960 frames per night.
\item \textbf{Label every 10 minutes} on the 5-point scale below. Record
      transitions to the minute.
\item Take a phone photo of the sky at each label as an independent
      adjudicator.
\item Note obstructions explicitly: tree, roofline, power line, aircraft,
      contrail, car headlights, dew on the lens.
\item Leave it running when you go inside. This campaign does not need
      supervision, only labels --- set a 10-minute timer and step out.
\end{enumerate}

\begin{center}
\small
\begin{tabular}{@{}clll@{}}
\toprule
\textbf{Code} & \textbf{Class} & \textbf{Criterion} & \textbf{Binary} \\
\midrule
0 & clear     & no cloud visible; Milky Way structure visible if moonless & 0 \\
1 & thin      & cirrus or haze; stars visible through it; halo around bright stars & 0 \\
2 & moderate  & distinct cloud over $\sim$25--60\,\% of the frame; some stars hidden & 1 \\
3 & heavy     & $>$60\,\% covered; only the brightest stars & 1 \\
4 & overcast  & no stars visible & 1 \\
\bottomrule
\end{tabular}
\end{center}

\textbf{Recorded.} Frames + a label timeline + moon altitude/illumination
(computable) + azimuth + humidity + phone photos.

\textbf{Abort criteria.} Rain --- bring the camera in. Dew on the lens ---
photograph it once deliberately as a hard negative, then dry it and note the
timestamp of the wipe.

\textbf{Acceptance.} $\geq$500 labelled frames per night. Over the campaign,
$\geq$800 frames in each of the three rarer classes.

\textbf{Class balance strategy.} After each night, tally per class and pick the
next night to fill the thinnest class. In this climate, class~1 (thin cirrus) is
usually the hardest to get deliberately and class~4 (overcast) the easiest and
least useful --- cap overcast at about 15\,\% of the set.

\subsection*{RB-6 --- Trailing ladder (R1 / R2)}
\addcontentsline{toc}{subsection}{RB-6 --- Trailing ladder}

\textbf{Purpose.} TED training data with \textbf{exactly computed} elongation
labels, exploiting the Z114's lack of tracking.

\textbf{Preconditions.} BP-4 --- the plate scale must be solved, not assumed,
because the label is computed in arcseconds and consumed in pixels.

\textbf{The label.} For an untracked exposure of $t$ seconds at declination
$\delta$, the trail length is
\[
\ell_{\mathrm{arcsec}} = 15.041 \, \cos\delta \, t ,
\qquad
\ell_{\mathrm{px}} = \ell_{\mathrm{arcsec}} / s ,
\]
with $s$ the solved plate scale. The expected ellipticity of a star of
seeing-limited FWHM $w$ (in px) trailed by $\ell_{\mathrm{px}}$ is
approximately $e = 1 - w/\sqrt{w^2 + \ell_{\mathrm{px}}^2}$, which is what the
classical estimator will report and what the CNN must learn to agree with in
the ambiguous band.

\begin{enumerate}
\item Choose \textbf{3 fields} at declinations near $0^\circ$, $+40^\circ$ and
      $+70^\circ$. $\cos\delta$ is 1.00, 0.77 and 0.34 --- so the same exposure
      gives three different trail lengths and the label grid is 2-D.
\item At each field, run the exposure ladder below. \textbf{Sample densely
      between 0.05 and 0.25\,s} --- that is where the CNN's ambiguous band
      $0.10 < e < 0.20$ falls, and a linear ladder under-samples it badly.
\begin{center}
\small
\begin{tabular}{@{}l@{}}
\toprule
$t$ (s) $\in \{0.017,\ 0.033,\ 0.05,\ 0.07,\ 0.10,\ 0.13,\ 0.17,\ 0.20,\ 0.25,\ 0.33,\ 0.50,\ 1.00\}$ \\
\bottomrule
\end{tabular}
\end{center}
\item \textbf{20 frames} at each exposure step. 3 fields $\times$ 12 steps
      $\times$ 20 = 720 frames per night.
\item Record for every frame: exposure, field declination, timestamp, and
      \textbf{tracking state = off}.
\item \textbf{Position angle control.} The trail direction is along the
      hour-angle axis, which for an alt-az tabletop mount rotates relative to
      the sensor as the target moves. Shoot the same field early and late in the
      evening to get two different position angles --- this prevents the model
      from learning ``elongation means horizontal''.
\item \textbf{Negatives.} Add either (a) $t \leq 0.02$\,s frames, where the trail
      is under half a pixel, or (b) frames from the tracked loaner rig. At least
      1,000 negative frames.
\item \textbf{Confounders (one session, deliberately):}
      \begin{itemize}
      \item Frames using the \emph{edge} of the field, where an f/3.95
            Newtonian's coma elongates stars radially --- a different position
            angle signature from tracking error.
      \item A windy night, which produces incoherent smear rather than a linear
            trail.
      \item A deliberate mid-exposure nudge (tap the tube), giving a double
            image rather than a streak.
      \end{itemize}
\end{enumerate}

\textbf{Recorded.} Frames + exposure + declination + tracking state + wind
estimate + which confounder class, if any.

\textbf{Abort criteria.} Cloud; star count $<10$ per frame.

\textbf{Acceptance.} Measured ellipticity versus predicted
$\ell_{\mathrm{px}}$ follows the expected curve within scatter. \textbf{If it
does not, the plate scale from BP-4 is wrong} --- and finding that out here is
much cheaper than finding it out during training.

\subsection*{RB-7 --- Sky background and lunation ladder}
\addcontentsline{toc}{subsection}{RB-7 --- Sky background ladder}

\textbf{Purpose.} Populate the sky-brightness covariate that both the exposure
bandit and the filter recommender consume, and validate the classical
background estimator against real gradients.

\textbf{Preconditions.} None beyond a body and tripod. Runs alongside RB-5 on
the same frames if the settings are held fixed.

\begin{enumerate}
\item Same fixed field, same exposure, same ISO, \textbf{once per hour, all
      night}.
\item Repeat across a full lunation: at least one night near new moon, one at
      first quarter, one near full, one at last quarter. Four nights spread over
      four weeks.
\item Include one session pointed low toward the light dome and one high toward
      zenith on the same night --- the gradient between them is the
      light-pollution signal.
\item Record moon altitude, illuminated fraction, and moon--field angular
      separation (all computable from time and site, but record the site
      coordinates).
\end{enumerate}

\textbf{Acceptance.} Background ADU varies monotonically with moon illumination
at fixed altitude. If it does not, something is auto-adjusting --- check that
exposure, ISO and white balance are truly locked.

\subsection*{RB-8 --- Exposure bandit logging (R5, rental)}
\addcontentsline{toc}{subsection}{RB-8 --- Bandit logging}

\textbf{Purpose.} Build a logged-bandit dataset that supports \emph{offline}
policy evaluation, so that a policy can be validated without spending rental
nights on live exploration.

\textbf{Preconditions.} Tracked mount; a body that tethers with bulb (T3i/T6S
are the safest bets after BP-1); D6 ideally resolved so airmass and moon
separation come from a solved pointing.

\begin{enumerate}
\item \textbf{Phase 1 --- uniform logging.} For the first 200 subs, select the
      exposure arm \textbf{uniformly at random} from
      $\{10, 30, 60, 90, 120, 180, 300\}$\,s. Record the arm \emph{and} its
      selection probability $1/7$.
\item \textbf{Phase 2 --- $\varepsilon$-greedy.} Thereafter use
      $\varepsilon = 0.2$: with probability 0.2 choose uniformly, otherwise
      choose the current best. \textbf{Record the realised probability for the
      chosen arm on every decision} --- this is the propensity, and without it
      no unbiased off-policy estimate exists.
\item \textbf{Record the 17-dim context} at decision time, before the exposure:
      sky background, moon illumination, $\cos\theta_{\mathrm{moon}}$, filter
      one-hot, target-type one-hot, seeing (star FWHM), airmass, guiding RMS.
      If there is no guider, record guiding RMS as \textbf{absent}, not zero.
\item \textbf{Record all five reward components separately} ---
      sub accepted, SNR in range, sub rejected, saturation $>5\%$, trailing ---
      in \co{frame\_measurements.payload\_json}. Never store only the scalar
      sum.
\item Vary target type across sessions --- galaxy, nebula, open cluster,
      globular --- since target type is four of the seventeen context
      dimensions and a single-target campaign leaves three of them constant.
\end{enumerate}

\textbf{Acceptance.} Every logged tuple carries a propensity. Every arm has
been selected at least 20 times. At least three distinct moon phases and two
target types are represented.

\textbf{Abort criteria.} If the propensity is not being logged, \textbf{stop and
fix it before continuing} --- data without it cannot be used for anything except
online learning, and the entire reason for the campaign is to avoid learning
online on rented time.

\subsection*{RB-9 --- Filter sequence (descoped)}
\addcontentsline{toc}{subsection}{RB-9 --- Filter sequence (descoped)}

Documented for completeness; not scheduled. If a loaner rig arrives with a
filter wheel, log every switch through \co{filter\_switch\_events} (schema
already present, migration 003) with: time, from-filter, to-filter, reason,
sky conditions at the time, and the realised SNR of the subs on each side. That
costs nothing and preserves the option. Do not purchase hardware for FOR --- see
\S\ref{m:for}.

\section*{Appendix C --- Observation log schema}
\addcontentsline{toc}{section}{Appendix C --- Observation log schema}

\subsection*{C.1 Session header (once per night)}

\begin{center}
\small
\begin{tabular}{@{}p{0.30\textwidth}p{0.16\textwidth}p{0.44\textwidth}@{}}
\toprule
\textbf{Field} & \textbf{Type} & \textbf{Note} \\
\midrule
\co{session\_id} & uuid & generated \\
\co{night\_id} & date & local date of \emph{sunset}, so one night is one ID across midnight \\
\co{rig\_id} & string & R1--R5, matching \S\ref{sec:rigs} \\
\co{site\_lat}, \co{site\_lon}, \co{site\_elev} & float & needed for airmass and moon geometry \\
\co{telescope}, \co{camera}, \co{barlow}, \co{focuser} & string & the exact optical train \\
\co{plate\_scale\_arcsec\_px} & float & from BP-4, per configuration \\
\co{saturation\_adu} & int & from BP-2, per body and ISO \\
\co{observer} & string & \\
\co{start\_time}, \co{end\_time} & ISO 8601 with offset & \\
\bottomrule
\end{tabular}
\end{center}

\subsection*{C.2 Environment series (every 15 minutes)}

\co{timestamp}, \co{ambient\_temp\_c}, \co{humidity\_pct}, \co{dew\_point\_c},
\co{wind\_beaufort}, \co{sky\_class} (0--4 per RB-5),
\co{transparency\_naked\_eye\_limiting\_mag}, \co{seeing\_impression}
(1--5, subjective, but recorded consistently), \co{notes}.

\begin{keybox}
Every one of these may be \textbf{absent}. Absent is a value; write
\co{null}, not a guess. \co{frames.temperature\_state} already models exactly
this distinction and the log must not undermine it.
\end{keybox}

\subsection*{C.3 Run record (per focus run or per campaign block)}

\co{run\_id}, \co{run\_type} (\co{focus\_sweep} $\vert$ \co{trailing\_ladder}
$\vert$ \co{cloud\_block} $\vert$ \co{drift} $\vert$ \co{bandit} $\vert$
\co{calibration}), \co{target\_name}, \co{target\_ra}, \co{target\_dec},
\co{target\_alt}, \co{target\_az}, \co{start\_position}, \co{step\_size},
\co{n\_positions}, \co{direction}, \co{tracking\_on}, \co{exposure\_s},
\co{iso\_or\_gain}, \co{filter}, \co{status}
(\co{completed} $\vert$ \co{aborted}), \co{abort\_reason}.

\subsection*{C.4 Disturbance record (as they happen)}

\co{timestamp}, \co{kind} (\co{bump} $\vert$ \co{refocus} $\vert$
\co{cloud} $\vert$ \co{headlights} $\vert$ \co{dew} $\vert$ \co{battery} $\vert$
\co{wind\_gust} $\vert$ \co{other}), \co{affected\_run\_id}, \co{description}.

\textbf{This is the block that gets skipped and should not be.} One unlogged
bump in the middle of a sweep produces a bimodal V-curve that the fitter will
resolve to the wrong minimum, and there is no way to detect it from the pixels
alone.

\subsection*{C.5 Machine-readable form}

The sidecar importer (\S G2/G4) consumes one JSON file per session:

\begin{lstlisting}[language=]
{
  "session": { "session_id": "...", "night_id": "2026-09-14",
               "rig_id": "R2", "site_lat": 36.07, "site_lon": -79.81,
               "telescope": "Zhumell Z114 + 2x Barlow",
               "camera": "SVBony SV105",
               "plate_scale_arcsec_px": 0.71,
               "saturation_adu": 255 },
  "environment": [
    { "t": "2026-09-14T21:15:00-04:00", "ambient_temp_c": 21.4,
      "humidity_pct": 68, "wind_beaufort": 1, "sky_class": 0,
      "seeing_impression": 3 },
    { "t": "2026-09-14T21:30:00-04:00", "ambient_temp_c": 20.9,
      "humidity_pct": 71, "wind_beaufort": 1, "sky_class": 1,
      "seeing_impression": 3 }
  ],
  "runs": [
    { "run_id": "...", "run_type": "focus_sweep", "target_name": "M39",
      "start_position": 1240, "step_size": 14, "n_positions": 21,
      "direction": "outward", "tracking_on": false,
      "exposure_s": 0.15, "iso_or_gain": 300, "status": "completed",
      "positions": [
        { "index": 0, "reading_um": 1240, "t": "...", "files": ["IMG_0001.CR2"] }
      ]
    }
  ],
  "disturbances": [
    { "t": "2026-09-14T23:02:00-04:00", "kind": "headlights",
      "affected_run_id": "...", "description": "car in driveway, 2 frames" }
  ]
}
\end{lstlisting}

\begin{keybox}
Fill the card on paper in the field and transcribe to JSON indoors the same
night. Do not plan to remember it in the morning --- you will not, and the
transcription step is also a useful review that catches missing entries while
the night is still recoverable.
\end{keybox}

\section*{Appendix D --- Post-session ingest checklist}
\addcontentsline{toc}{section}{Appendix D --- Post-session ingest checklist}

\begin{enumerate}
\item \textbf{Do not format the card.} Copy files off it; keep the originals
      until the checksums verify.
\item Transcribe the paper card to the JSON schema of \S C.5 \textbf{the same
      night}.
\item Copy frames to the capture root, preserving the camera's filenames.
\item Run the sidecar importer for anything the daemon did not capture itself.
      Confirm it stamps \co{driver\_name = "offline\_import"}.
\item Start the daemon and let \co{reconcile()} run. Read its log --- it reports
      committed rows with missing files and staged files with no row.
\item Verify checksums against \co{frame\_assets.checksum\_sha256}. A mismatch
      means a bad copy; recopy from the card, which is why step 1 exists.
\item Count what you got, against what you attempted:
\begin{lstlisting}[language=bash]
sqlite3 achernar.db "
  SELECT capture_type, capture_status, COUNT(*)
  FROM frames WHERE date(captured_at)='2026-09-14'
  GROUP BY 1,2;"
sqlite3 achernar.db "
  SELECT status, COUNT(*) FROM capture_attempts
  WHERE date(started_at)='2026-09-14' GROUP BY 1;"
\end{lstlisting}
\item Dry-run the dataset builder and \textbf{check the group count, not the row
      count}:
\begin{lstlisting}[language=bash]
python -m ml.data_pipeline.build_focus_quality_dataset \
  --db achernar.db --output /tmp/check.npz
python3 -c "import numpy as np; d=np.load('/tmp/check.npz',allow_pickle=True); \
  print('rows', d['crops'].shape[0], 'groups', len(set(d['run_key'].tolist())))"
\end{lstlisting}
\item Confirm the saturation gate actually ran on this night's frames. If it
      reports \emph{not performed}, D1 or D2 is still open and the night's focus
      labels are suspect.
\item Back up: capture root + \co{achernar.db} + the session JSON, to a second
      physical drive. Only then reuse the card.
\item Update the running tally --- runs per rig, frames per COC class, nights
      per campaign --- against the budgets in Part~III, and pick the next
      night's target accordingly.
\end{enumerate}

\section*{Appendix E --- Standing decisions}
\addcontentsline{toc}{section}{Appendix E --- Standing decisions}

\begin{center}
\small
\begin{tabular}{@{}p{0.30\textwidth}p{0.60\textwidth}@{}}
\toprule
\textbf{Decision} & \textbf{Rationale} \\
\midrule
Collect COC labels as 5-class & Collapsible to the spec's binary head; the reverse is impossible, and the sky is not re-observable (D7). \\
2$\times$ Barlow is mandatory on the Z114 & Native focal length is undersampled at best focus for every available camera (\S\ref{sec:sampling}). \\
Descope FOR (Model 8) & No wheel, no filters, \co{hardware/filter\_wheel.rs} is \co{todo!()}, and its features depend on models that do not exist yet (\S\ref{m:for}). \\
TME edge by distillation, not contrastive & AstroM3's published weights already exist; only the image adapter is untrained, and $10^5$ contrastive objects are out of reach (\S\ref{m:tme}). \\
FOP data by passive harvesting & 20 rows per run cannot reach $10^4$ from dedicated sweeps alone (\S\ref{m:fop}). \\
Ambient temperature logged by hand for v1 & \co{hardware/sensors.rs} is \co{todo!()}; manual logging is well within the Kalman filter's assumed $\sigma_z = 3$ steps (\S\ref{m:fdf}). \\
Dial indicator instead of a motorised focuser for v0/v1 & \$15 buys 10\,\textmu m absolute position, which makes manual runs FOP-eligible (RB-2). \\
Rental window spent last & It is the scarce resource; field procedures should be debugged on free gear first (\S\ref{sec:critical}). \\
\bottomrule
\end{tabular}
\end{center}

\vspace{1em}
\begin{keybox}
\textbf{Open questions for you to resolve, in priority order.}
\begin{enumerate}
\item Does the loan programme's rig include an electronic focuser, and on what
      protocol? This decides whether \$200--300 needs spending and whether T2
      and T3 merge.
\item Does \co{cargo build} pass on the Mac (BP-0)? Everything downstream is
      contingent.
\item Can libuvc open the SV105 on macOS (BP-1)? If not, R1 and R2 need a Linux
      host.
\item Which of the five bodies actually support remote capture (BP-1)?
\item Does any camera reach prime focus on the Z114 (BP-3)?
\end{enumerate}
\end{keybox}

\end{document}
